%%% Copyright Sheldon Axler, 2024. This file may be used only with the permission of Sheldon Axler.

\chapter{Inner Product Spaces} \label{89}

In making the definition of a vector space, we generalized the linear structure
(addition and scalar multiplication) of $\R{2}$ and $\RR{3}$.  We ignored geometric features such as the notions of length and angle.  These ideas are embedded in the concept of inner products, which we will investigate in this chapter.

Every inner product induces a norm, which you can think of as a length. This norm satisfies key properties such as the Pythagorean theorem, the triangle inequality, the parallelogram equality, and the Cauchy--Schwarz inequality.

The notion of perpendicular vectors in Euclidean geometry gets renamed to orthogonal vectors in the context of an inner product space. We will see that orthonormal bases are tremendously useful in inner product spaces. The Gram--Schmidt procedure constructs such bases. This chapter will conclude by putting together these tools to solve minimization problems.

\StandingBox{
\item
$\F{}$ denotes $\R{}$ or $\C{}$.
\item
$V$ and $W$ denote vector spaces over $\F{}\3$.
}

\begin{comment}
\TopicsBox{
\item
Cauchy--Schwarz inequality
\item
Gram--Schmidt procedure
\item
linear functionals on inner product spaces
\item
calculating minimum distance to a subspace
}
\end{comment}

\renewcommand{\FigureName}{\ifnum\sizepaper=1 Peabody.png\else PeabodyR.png \fi}
\renewcommand{\FigureScale}{\ifnum\sizepaper=1 0.114\else 1\fi}

\FigureHereCredit{artwork/\FigureName}{\FigureScale}%
{The George Peabody Library, now part of Johns Hopkins University, opened while James Sylvester \uppar{1814--1897} was the university's first mathematics professor. Sylvester's publications include the first use of the word \marbf{matrix} in mathematics.\index{Sylvester, James}\index{Peabody Library}}{1in}{Matthew Petroff CC BY-SA}

%\FigureHereCredit{artwork/WomanTeachingGeometry.png}{0.18}{Woman teaching geometry, from a fourteenth-century\\edition of Euclid's geometry book.}{1.1in}{British Library}

\section{Inner Products and Norms} \label{6ipn}

\subsection{Inner Products}

\marF[9]{\includegraphics{Mathematica/Axes.pdf}}{This vector\/ $v$ has norm\/ $ \sqrt{a^2 + b^2}$.}

To motivate the concept of inner product, think of vectors in~$\R{2}$
and~$\R{3}$ as arrows with initial point at the origin.  The length of a vector $v$
in~$\R{2}$ or~$\R{3}$ is called the \emph{norm} of~$v$ and is denoted by $\|v\|$.  Thus for
\mbox{$v = (a, b) \in \RR{2}$}, we have
\[
\|v\| =\sqrt{a^2 + b^2}.
\]
Similarly, if $v = (a,b,c) \in \RR{3}$, then $\|v\| = \sqrt{a^2 + b^2 + c^2}$.

Even though we cannot draw pictures in higher dimensions, the generalization to~$\R{n}$ is easy: we define the norm of
$x = (x_1, \dots, x_n) \in \R{n}$ by\index{norm of a vector}
\[
\|x\| = \sqrt{{x_1}^{\2} + \dots + {x_n}^{\2}}.
\]

The norm is not linear on $\RR{n}$.  To inject linearity into the discussion, we
introduce the dot product.

\DefinitionBox{dot product}{For $x, y \in \RR{n}$, the
\emph{dot product}\index{dot product} of~$x$ and~$y$, denoted by $x \cdot y$, is
defined by
\[
x \cdot y = x_1 y_1 + \dots + x_n y_n,
\]
where $x = (x_1, \dots, x_n)$ and $y = (y_1, \dots, y_n)$.}

\mar{If we think of a vector as a point \mbox{instead} of as an arrow, then\/ $\|x\|$ should be \mbox{interpreted} to mean the distance from the origin to the point\/ $x$.}

The dot product of two vectors in $\R{n}$ is a number, not a vector. Notice that $x \cdot x = \|x\|^2$ for all $x \in \RR{n}$.  Furthermore, the dot product on $\R{n}$ has the following properties.
\begin{Itemize}
\item
$x \cdot x \ge 0$ for all $x \in \RR{n}$.
\item
$x \cdot x = 0$ if and only if $x = 0$.
\item
For $y \in \R{n}$ fixed, the map from~$\R{n}$ to~$\R{}$ that sends $x \in \R{n}$ to $x \cdot y$ is linear.
\item
$x \cdot y = y \cdot x$ for all $x, y \in \RR{n}$.
\end{Itemize}

An inner product is a generalization of the dot product.  At this point you may
be tempted to guess that an inner product is defined by abstracting the
properties of the dot product discussed in the last paragraph.  For real vector
spaces, that guess is correct.  However, so that we can make a definition that will
be useful for both real and complex vector spaces, we need to examine the
complex case before making the definition.

\pagebreak[1]

Recall that if $\la = a + bi$, where $a, b \in \R{}$, then
\begin{Itemize}
\item
the absolute value of $\la$, denoted by $|\la|$, is defined by $|\la| = \sqrt{a^2 + b^2}$;
\item
the complex conjugate of $\la$, denoted by $\bar{\la}$, is defined by $\bar{\la} = a - bi$;
\item
$|\la|^2 = \la \bar{\la}$.
\end{Itemize}
See Chapter~\ref{173} for the definitions and the
basic properties of the absolute value and complex conjugate.

For
\mbox{$z = (z_1, \dots, z_n) \in \CC{n}$}, we define the norm of $z$ by
\[
\|z\| = \sqrt{|z_1|^2 + \dots + |z_n|^2}.
\]
The absolute values are needed because we want $\|z\|$ to be a nonnegative
number.  Note that
\[
\|z\|^2 = z_1 \overline{z_1} + \dots + z_n \overline{z_n}.
\]

We want to think of $\|z\|^2$ as the inner product of $z$ with itself, as we did
in~$\RR{n}$.  The equation above thus suggests that the inner product of the vector
\mbox{$w = (w_1, \dots, w_n) \in \C{n}$} with $z$ should equal
\[
w_1 \overline{z_1}+ \dots + w_n \overline{z_n}.
\]
If the roles of the $w$ and $z$ were interchanged, the expression above would be
replaced with its complex conjugate. Thus we should expect that the
inner product of~$w$ with~$z$ equals the complex conjugate of the inner product
of~$z$ with~$w$.  With that motivation, we are now ready to define an inner
product on~$V\1$, which may be a real or a complex vector space.

One comment about the notation used in the next definition:
\begin{Itemize}
\item
For $\la\in\C{}$, the notation $\la \ge 0$ means $\la$ is real and~nonnegative.
%\item
%We use the common notation $\ip{u, v}$, with angle brackets denoting an inner product. Some people use parentheses instead, but %then $(u, v)$ becomes ambiguous because it could denote either an ordered pair or an inner product.
\end{Itemize}

\DefinitionBox[innerproduct]{inner product}{An \emph{inner product}\index{inner product} on $V$ is a function that takes each ordered pair $(u,v)$ of  elements of $V$ to a number $\ip{u,v} \in \F{}$ and has the following properties.
\begin{description} \addtolength{\itemsep}{-1bp}
\item[positivity] \ \\
$\ip{v, v} \ge 0$ for all $v \in V\1$.

\item[definiteness] \ \\
$\ip{v, v} = 0$ if and only if $v = 0$.

\item[additivity in first slot] \ \\
$\ip{u + v, w} = \ip{u, w} + \ip{v, w}$ for all $u, v, w \in V\1$.

\item[homogeneity in first slot] \ \\
$\ip{\la u, v} = \la \ip{u, v}$ for all $\la \in \F{}$ and all $u, v \in V\1$.

\item[conjugate symmetry] \ \\
$\ip{u, v} = \overline{\ip{v,u}}$ for all $u, v \in V\1$.\index{conjugate symmetry}
\end{description}
}\sindex{uv@$\ip{u, v}$}

\mar{Most mathematicians define inner products as above, but many physicists use a definition that requires homogeneity in the second slot instead of the first slot.}

Every real number equals its complex conjugate.  Thus if we are
dealing with a real vector space, then in the  last condition above we can dispense
with the complex conjugate and simply state that $\ip{u,v} = \ip{v,u}$ for all
$u, v \in V\1$.

\begin{example}[exampleip]{inner products}
\begin{subtheorem}*
\item
The \emph{Euclidean inner product}\index{Euclidean inner product} on~$\F{n}$ is defined by
\[
\ip[\big]{(w_1, \dots, w_n), (z_1, \dots, z_n)} =
w_1 \overline{z_1}+ \dots + w_n \overline{z_n}
\]
for all $(w_1, \dots, w_n), (z_1, \dots, z_n) \in \FF n$.

\item
If $c_1, \dots, c_n$ are positive
numbers, then an inner product can be defined on~$\F{n}$ by
\[
\ip[\big]{(w_1, \dots, w_n), (z_1, \dots, z_n)} =
c_1 w_1 \overline{z_1} + \dots + c_n w_n \overline{z_n}
\]
for all $(w_1, \dots, w_n), (z_1, \dots, z_n) \in \FF n$.

\item
An inner product can be defined on the vector space of continuous real-valued functions on the interval $[-1, 1]$ by
\[
\ip{f, g} = \int_{-1}^1fg
\]
for all $f, g$ continuous real-valued functions on $[-1,1]$.

\item
An inner product can be defined on $\P(\R{})$ by
\[
\ip{p, q} = p(0) q(0) + \int_{-1}^1 p' q'
\]
for all $p, q \in \P(\R{})$.

\item
An inner product can be defined on $\P(\R{})$ by
\[
\ip{p, q} = \int_0^\infty p(x) q(x) e^{-x} \, dx
\]
for all $p, q \in \P(\R{})$.
\end{subtheorem}
\exampleitemend
\end{example}

\DefinitionBox{inner product space}{An \emph{inner product space}\index{inner product space} is a vector space $V$ along with an inner product on $V\1$.}

The most important example of an inner product space is $\F{n}$ with the Euclidean inner product given by (a) in the example above. When $\F{n}$ is referred to as an inner product space, you should assume that the inner product is the Euclidean inner product unless explicitly told otherwise.

\pagebreak[1]

So that we do not have to keep repeating the hypothesis that $V$ and $W$ are inner product spaces, we make the following assumption.

\NotationBox{$V\1$, $W$}{For the rest of this chapter and the next chapter, $V$ and $W$ denote inner product spaces over $\F{}\3$.}

Note the slight abuse of language here. An inner product space is a vector space along with an inner product on that vector space. When we say that a vector space $V$ is an inner product space, we are also thinking that an inner product on $V$ is lurking nearby or is clear from the context \big(or is the Euclidean inner product if the vector space is $\F{n}$\big).

\ResultBox[basicip]{basic properties of an inner product}
{\begin{subtheorem}
\item
For each fixed $v \in V\1$, the function that takes $u \in V$ to $\ip{u,v}$ is a linear map from $V$ to $\F{}\3$.
\item
$\ip{0, v} = 0$ for every $v \in V\1$.
\item
$\ip{v, 0} = 0$ for every $v \in V\1$.
\item
$\ip{u, v + w} = \ip{u , v} + \ip{u, w}$ for all $u, v, w \in V\1$.
\item
$\ip{u, \la v} = \bar{\la} \ip{u,v}$ for all $\la \in \F{}$ and all $u, v \in V\1$.
\end{subtheorem}}

\proof
\begin{subtheorem}*
\item
For $v \in V\1$, the linearity of $u \mapsto \ip{u,v}$ follows from the conditions of additivity and homogeneity in the first slot in the definition of an inner product.

\item
Every linear map takes $0$ to $0$. Thus (b) follows from (a).

\item
If $v \in V\1$, then the conjugate symmetry property in the definition of an inner product and (b) show that $\ip{v, 0} = \overline{ \ip{0, v} } = \overline{0} = 0$.

\item
Suppose $u, v, w \in V\1$. Then
\begin{align*}
\ip{u, v + w} &= \overline{\ip{v + w, u}} \\[2bp]
&= \overline{\ip{v , u} + \ip{w, u}} \\[2bp]
&= \overline{\ip{v , u}} + \overline{\ip{w, u}} \\[2bp]
&= \ip{u , v} + \ip{u, w}.
\end{align*}

\item
Suppose $\la \in \F{}$ and $u,v \in V\1$. Then
\begin{align*}
\ip{u, \la v} &= \overline{\ip{\la v,u}} \\[2bp]
&= \overline{\la\ip{v,u}} \\[2bp]
&= \bar{\la} \, \, \overline{\ip{v,u}} \\[2bp]
&= \bar{\la} \ip{u,v}. \rlap{ \hspace{1.5in} \qedsymbolMath} %VisualFormatting for LADR4e
\end{align*}
\end{subtheorem}

\pagebreak[3]

\subsection{Norms}

Our motivation for defining inner products came initially from the norms of vectors on $\R2$ and $\RR3$. Now we see that each inner product determines a norm.

\DefinitionBox{norm, $\|v\|$}
{For $v \in V\1$, the \emph{norm}\index{norm of a vector} of $v$, denoted by $\|v\|$, is defined by
\[
\|v\| = \sqrt{\ip{v,v}}.
\]}\sindex{w@$"|\hspace{-1bp}"|v"|\hspace{-1bp}"|$}

\begin{example}{norms}
\begin{subtheorem}*
\item
If $(z_1, \dots, z_n) \in \F{n}$ (with the Euclidean inner
product), then
\[
\|(z_1, \dots, z_n)\| = \sqrt{|z_1|^2 + \dots + |z_n|^2}.
\]

\item
For $f$ in the vector space of continuous real-valued functions on $[-1,1]$ and with inner product given as in \ref{exampleip}(c), we have
\[
\|f\| =  \sqrt{\int_{-1}^1 f^2}.
\]
\end{subtheorem}
\vspace{-6bp}
\end{example}

\ResultBox[2bnorm]{basic properties of the norm}
{Suppose $v \in V\1$.
\begin{subtheorem}*
\item
$\|v\| = 0$ if and only if $v = 0$.
\item
$\|\la v\| = |\la|\,\|v\|$ for all $\la \in \F{}\3$.
\end{subtheorem}}

\proof
\begin{subtheorem}*
\item
The desired result holds because $\ip{v,v} = 0$ if and only if $v = 0$.

\item
Suppose $\la \in \F{}\3$. Then
\begin{align*}
\|\la v\|^2 &= \ip{\la v, \la v} \\[4bp]
&= \la \ip{v, \la v} \\[4bp]
&= \la \bar{\la}\ip{v, v} \\[4bp]
&= |\la|^2 \, \|v\|^2\1.
\end{align*}
Taking square roots now gives the desired equality.
\qed
\end{subtheorem}

The proof of (b) in the result above illustrates a general principle: working with norms squared is usually easier than working directly with norms.

\pagebreak[3]

Now we come to a crucial definition.

\DefinitionBox{orthogonal}{Two vectors $u, v \in V$ are called
\emph{orthogonal}\index{orthogonal!vectors} if $\ip{u, v} = 0$.}

\mar[5]{The word \marbf{orthogonal} comes from the Greek word \marbf{orthogonios},
which means right-angled.}

In the definition above, the order of the two vectors does not matter, because $\ip{u, v} = 0$ if and only if $\ip{v, u} = 0$. Instead of saying $u$ and $v$ are orthogonal, sometimes we say $u$ is  orthogonal to $v$.

Exercise \ref{cost} asks you to prove that if $u, v$ are nonzero vectors in $\RR{2}$, then
\[
\ip{u,v} = \|u\| \, \|v\| \cos \theta,
\]
where $\theta$ is the angle between $u$ and $v$ (thinking of~$u$ and~$v$ as arrows
with initial point at the origin). Thus two nonzero vectors in $\R2$ are orthogonal (with respect to the Euclidean inner product) if and only if the cosine of the angle between them is $0$, which happens if and only if the vectors are perpendicular in the usual sense of plane geometry. Thus you can think of the word \textit{orthogonal} as a fancy word meaning \textit{perpendicular}.

We begin our study of orthogonality with an easy result.

\ResultBox{orthogonality and $0$}
{\begin{subtheorem}
\item
$0$ is orthogonal to every vector in $V\1$.
\item
$0$ is the only vector in $V$ that is orthogonal to itself.
\end{subtheorem}}

\proof
\begin{subtheorem}*
\item
Recall that \ref{basicip}(b) states that $\ip{0, v} = 0$ for every $v \in V\1$.
\item
If $v \in V$ and $\ip{v,v} = 0$, then $v = 0$ (by definition of inner product).
\qed
\end{subtheorem}

For the special case $V = \RR{2}$, the next theorem was known over 3,500 years ago in Babylonia and then rediscovered and proved over 2,500 years ago in Greece. Of course, the proof below is not the original proof.

\ResultBox[184]{Pythagorean theorem}{\index{Pythagorean theorem}%
Suppose $u, v \in V\1$. If $u$ and $v$ are orthogonal, then
\[
\|u + v\|^2 = \|u\|^2 + \|v\|^2\1.
\]
}

\proof
Suppose $\ip{u, v} = 0$. Then
\begin{align*}
\|u + v\|^2 &= \ip{u+ v, u + v} \\
&= \ip{u, u} + \ip{u, v} + \ip{v, u} + \ip{v, v} \\
&= \|u\|^2 + \|v\|^2\1. \rlap{ \hspace{2.03in} \qedsymbolMath} %VisualFormatting for LADR4e
\end{align*}

\pagebreak[3]\par\addvspace{\listsep} %To go along with \qedsymbolMath above.

Suppose $u, v \in V\1$, with $v \ne 0$.  We would like to write $u$ as a scalar multiple of~$v$ plus a
vector $w$ orthogonal to~$v$, as suggested in the picture here.

\FigureHereUnround{Mathematica/OrthogonalDecomposition.pdf}{1}{An orthogonal decomposition:\\
$u$ expressed as a scalar multiple of\/ $v$ plus a vector orthogonal to\/ $v$.}

To discover how to write $u$ as a scalar multiple of~$v$ plus a vector orthogonal
to~$v$, let $c \in \F{}$ denote a scalar.  Then
\[
u = c v + (u - c v).
\]
Thus we need to choose $c$ so that $v$ is orthogonal to $(u - c v)$. Hence we want
\[
0 = \ip{u - c v, v} = \ip{u,v} - c \|v\|^2\1.
\]
The equation above shows that we should choose $c$ to be $\ip{u,v}/\|v\|^2\1$.  Making this choice of~$c$, we
can write
\[
u = \frac{\ip{u,v}}{\|v\|^2} v + \paren[\bigg]{ u - \frac{\ip{u,v}}{\|v\|^2} v }.
\]
As you should verify, the equation displayed above explicitly writes $u$ as a scalar
multiple of~$v$ plus a vector orthogonal to~$v$. Thus we have proved the following key result.

\ResultBox[orthogdecomp]{an orthogonal decomposition}
{Suppose $u, v \in V\1$, with $v \ne 0$. Set $c =  \dfrac{\ip{u,v}}{\|v\|^2}$ and $w = u - \dfrac{\ip{u,v}}{\|v\|^2} v$. Then
\[
u = c v + w \andd \ip{w, v} = 0.
\]}

The orthogonal decomposition \ref{orthogdecomp} will be used in the proof of the Cauchy--Schwarz inequality, which is our next result and is one of the most important inequalities in mathematics.

\ResultBox[177]{Cauchy--Schwarz inequality}{\index{Cauchy--Schwarz inequality}%
Suppose $u, v \in V\1$. Then
\[
|\ip{u, v}| \le \|u\|\,\|v\|.
\]
This inequality is an equality if and only if one of\/ $u, v$ is a
scalar multiple of the other.
}

\proof
If $v = 0$, then both sides of the desired inequality equal $0$.  Thus we can
assume that $v \ne 0$.  Consider the orthogonal decomposition
\[
u = \frac{\ip{u,v}}{\|v\|^2} v + w
\]
given by \ref{orthogdecomp}, where $w$ is orthogonal to $v$.  By the Pythagorean theorem,
\begin{align}
\|u\|^2 &= \norm[\bigg]{ \frac{\ip{u,v}}{\|v\|^2} v }^2 + \norm{w}^2 \notag \\[6bp]
&= \frac{\abs[\big]{\ip{u,v}}^2}{\|v\|^2}  + \|w\|^2 \notag \\[6bp]
&\ge \frac{\abs[\big]{\ip{u,v}}^2}{\|v\|^2}. \label{336}
\end{align}
Multiplying both sides of this inequality by $\|v\|^2$ and  then taking square roots gives the desired inequality.

\mar{Augustin-Louis Cauchy \uppar{1789--1857}
proved \ref{CS}\uppar{a} in 1821. In 1859, Cauchy's student Viktor Bunyakovsky \uppar{1804--1889} proved integral inequalities like the one in \ref{CS}\uppar{b}. A few decades later, similar discoveries by Hermann Schwarz \uppar{1843--1921} attracted more attention and led to the name of this inequality.\index{Cauchy, Augustin-Louis}\index{Bunyakovsky, Viktor}\index{Schwarz, Hermann}}

The proof in the paragraph above shows that the Cauchy--Schwarz inequality is an equality if and only if \ref{336} is an equality. This happens if and only if $w = 0$.  But $w=0$ if and only if $u$ is a multiple of $v$ (see~\ref{orthogdecomp}).  Thus the
Cauchy--Schwarz inequality is an equality if and only if $u$ is a scalar multiple of
$v$ or $v$ is a scalar multiple of $u$ (or both; the phrasing has been chosen to
cover cases in which either $u$ or $v$ equals~$0$).
\qed

\begin{example}[CS]{Cauchy--Schwarz inequality}
\begin{subtheorem}*
\item
If $x_1, \ldots, x_n, y_1, \ldots, y_n \in \R{}$, then
\[
( x_1 y_1 + \cdots + x_n y_n)^2 \le \paren[\big]{ {x_1}^{\2} + \cdots + {x_n}^{\2} } \paren[\big]{ {y_1}^{\2} + \cdots + {y_n}^{\2} },
\]
as follows from applying the Cauchy--Schwarz inequality to the vectors $(x_1, \ldots, x_n), (y_1, \ldots, y_n) \in \RR n$, using the usual Euclidean inner product.
\item
If $f, g$ are continuous real-valued functions on $[-1, 1]$, then
\[
\abs[\Big]{\int_{-1}^1 fg \,}^{\,2}
\le \paren[\Big]{ \int_{-1}^1 f^2 }
\paren[\Big]{ \int_{-1}^1 g^2 },
\]
as follows from applying the Cauchy--Schwarz inequality to Example \ref{exampleip}(c).
\end{subtheorem}
\vspace{-8bp}
\end{example}

\begin{comment}
\begin{example}[CS]{Cauchy--Schwarz inequality}
\begin{subtheorem}
\item
If $x_1, \dots, x_n, y_1, \dots, y_n \in \R{}$, then
\[
|x_1 y_1 + \dots + x_n y_n|^2 \le
\paren[\big]{{x_1}^{\2} + \dots + {x_n}^{\2}}\paren[\big]{{y_1}^{\2} + \dots + {y_n}^{\2}},
\]
as follows from applying the Cauchy--Schwarz inequality to $\R{n}$ with the Euclidean inner product.
\item
If $f, g$ are continuous real-valued functions on $[-1, 1]$, then
\[
\abs[\bigg]{\int_{-1}^1 fg }^2
\le \paren[\bigg]{ \int_{-1}^1 f^2  }
\paren[\bigg]{ \int_{-1}^1 \g^2  },
\]
as follows from applying the Cauchy--Schwarz inequality to the inner product space in Example \ref{exampleip}(c).
\end{subtheorem}
\end{example}
\end{comment}

\marF[7]{\includegraphics{Mathematica/TriangleInequality3.pdf}}{\vspace*{-.2in}In this triangle, the length of\\$u+v$ is less than the length\\of\/ $u$ plus the length of\/ $v$.}

The next result, called the triangle inequality\index{triangle inequality}, has the geometric interpretation that the length of each side of a triangle is less than the sum of the lengths of the other two sides.

Note that the triangle inequality implies that the shortest polygonal path between two points is a single line segment (a polygonal path consists of line segments).

\vspace{1bp}

\ResultBox[3triangle]{triangle inequality}{
Suppose $u, v \in V\1$. Then
\vspace{-.06in}
\[
\|u + v\| \le \|u\| + \|v\|.
\]
This inequality is an equality if and only if one of\/ $u, v$ is a nonnegative real multiple of the other.
}

\proof
We have
\vspace{-3bp}
\begin{align}
\|u + v\|^2 &= \ip{u + v, u + v} \notag \\[1bp]
&=\ip{u, u} + \ip{v, v} + \ip{u, v} + \ip{v, u} \notag \\[1bp]
&=\ip{u, u} + \ip{v, v} + \ip{u, v} + \overline{\ip{u, v}} \notag \\[1bp]
&= \|u\|^2 + \|v\|^2 + 2\re \ip{u, v} \notag \\[1bp]
&\le \|u\|^2 + \|v\|^2 + 2 \abs[\big]{ \ip{u, v} } \label{176} \\[1bp]
&\le \|u\|^2 + \|v\|^2 + 2 \|u\|\,\|v\| \label{42} \\[1bp]
&= \paren[\big]{\|u\| + \|v\|}^2\1,\notag
\end{align}
where \ref{42} follows from the Cauchy--Schwarz inequality (\ref{177}).  Taking
square roots of both sides of the inequality above gives the desired inequality.

The proof above shows that the triangle inequality is an equality if
and only if we have equality in \ref{176} and~\ref{42}.  Thus we have equality
in the triangle inequality if and only if
\begin{equation} \label{374}
\ip{u, v} = \|u\| \, \|v\|.
\end{equation}%
If one of $u, v$ is a nonnegative real multiple of the other, then \ref{374} holds.  Conversely, suppose \ref{374} holds.  Then the condition for equality in the Cauchy--Schwarz inequality (\ref{177}) implies that one of $u, v$
is a scalar multiple of the other. This scalar must be a nonnegative real number, by \ref{374}, completing the proof.
\qed

For the reverse triangle inequality, see Exercise \ref{6reverse}.

\pagebreak[3]

\marF[7]{\includegraphics{Mathematica/Parallelogram.pdf}}{The diagonals of this parallelogram\\are\/ $u+v$ and\/ $u-v$.}

The next result is called the parallelogram equality because of its geometric interpretation: in every parallelogram, the sum of the squares of the lengths of the diagonals equals the sum of the squares of the lengths of the four sides. Note that the proof here is more straightforward than the usual proof in Euclidean geometry.

\ResultBox{parallelogram equality}{
Suppose $u, v \in V\1$. Then\index{parallelogram equality}
\[
\|u + v\|^2 + \|u - v\|^2 = 2(\|u\|^2 + \|v\|^2).
\]
}

\proof
We have
\begin{align*}
\|u + v\|^2 + \|u - v\|^2 &= \ip{u + v, u + v} + \ip{u - v, u - v} \\[5bp]
&= \|u\|^2 + \|v\|^2 + \ip{u, v} + \ip{v,u} \\
&\quad \quad + \|u\|^2 + \|v\|^2 - \ip{u,  v} -\ip{v,u} \\[5bp]
&= 2 \paren[\big]{\|u\|^2 + \|v\|^2},
\end{align*}
as desired.
\qed

\begin{exercises}

\exercise
Prove or give a counterexample: If $v_1, \dots, v_m \in V\1$, then
\[
\sum_{\js}^m \sum_{k=1}^m \ip{v_j, v_k} \ge 0.
\]
\vspace{-12bp}


\exercise
Suppose $S \in \L(V)$. Define $\ip{\cdot, \cdot}_1$ by
\[
\ip{u, v}_1 = \ip{Su, Sv}
\]
for all $u, v \in V\1$. Show that $\ip{\cdot, \cdot}_1$ is an inner product on $V$ if and only if $S$ is injective.


\exercise
\begin{SUBtheorem}
\item
Show that the function taking an ordered pair $\bigl( (x_1, x_2), (y_1, y_2) \bigr)$ of elements of $\R2$ to $|x_1 y_1| + |x_2 y_2|$ is not an inner product on $\RR2$.
\item
Show that the function taking an ordered pair $\bigl( (x_1, x_2, x_3), (y_1, y_2, y_3) \bigr)$ of elements of $\R3$ to $x_1 y_1 + x_3 y_3$ is not an inner product on $\RR3$.
\end{SUBtheorem}


\exercise
Suppose $T \in \L(V)$ is such that $\|Tv\| \le \|v\|$ for every $v \in V\1$. Prove that $ T - \sqrt{2}\, I$ is injective.


\exercise
Suppose $V$ is a real inner product space.
\begin{subtheorem}
\item
Show that $\ip{u + v, u - v} = \| u \|^2 - \| v \|^2$ for every $u, v \in V\1$.
\item
Show that if $u, v \in V$ have the same norm, then $u+v$ is orthogonal to $u-v$.
\item
Use (b) to show that the diagonals of a rhombus are perpendicular to each other.
\end{subtheorem}
\exercisesubtheoremend


\exercise
Suppose $u, v \in V\1$.  Prove that $\ip{u,v} = 0 \iff
\|u\| \le \|u + av\|
$
for all $a \in \F{}\3$.\label{orthognorm}


\exercise
Suppose $u, v \in V\1$. Prove that $\|au + bv\| = \|bu + av\|$ for all $a, b \in \R{}$ if and only if $\|u\| = \|v\|$.


\exercise
Suppose $a,b,c,x, y \in \R{}$ and $a^2 + b^2 + c^2 + x^2 + y^2 \le 1$. Prove that $a + b + c +4x + 9y \le 10$.


\exercise
Suppose $u, v \in V$ and $\| u \| = \|v \| = 1$ and $\ip{u, v} = 1$. Prove that $u = v$.


\exercise
Suppose $u, v \in V$ and $\|u\| \le 1$ and $\|v\| \le 1$. Prove that
\[
\sqrt{ 1 - \|u\|^2} \sqrt{ 1 - \|v\|^2} \le 1 - \abs[\big]{\ip{u, v}}.
\]
\exercisedisplayend


\exercise
Find vectors $u, v \in \R{2}$ such that $u$ is a scalar multiple of $(1,3)$, $v$ is orthogonal to $(1,3)$, and $(1,2) = u + v$.


\exercise
Suppose $a, b, c, d$ are positive numbers.
\begin{subtheorem}
\item
Prove that
$ \displaystyle
(a + b + c + d)\paren[\bigg]{\frac{1}{a} + \frac{1}{b} + \frac{1}{c} + \frac{1}{d}} \ge 16$.
\item
For which positive numbers $a, b, c, d$ is the inequality above an equality?
\end{subtheorem}
\exercisesubtheoremend


\exercise
Show that the square of an average is less than or equal to the average of the squares. More precisely, show that if $a_1, \dots , a_n \in \R{}$, then the square of the average of $a_1, \dots , a_n$ is less than or equal to the average of ${a_1}^{\2}\1, \dots, {a_n}^{\2}\1$.


\exercise
Suppose $v \in V$ and $v \ne 0$. Prove that $v / \norm{v}$ is the unique closest element on the unit sphere of $V$ to $v$. More precisely, prove that if $u \in V$ and $\norm{u} = 1$, then
\[
\norm[\bigg]{ v - \frac{v}{\norm{v}} } \le \norm{v - u},
\]
with equality only if $u = v / \norm{v}$.


\exercise
Suppose $u, v$ are nonzero vectors in $\RR{2}$. Prove that \label{cost}
\[
\ip{u, v} = \|u\| \, \|v\| \cos \theta,
\]
where $\theta$ is the angle between $u$ and $v$ (thinking of~$u$ and~$v$ as arrows
with initial point at the origin).
\hint{Use the law of cosines on the triangle formed by $u$, $v$, and~$u-v$.}


\exercise
The angle between two vectors (thought of as arrows with initial point at the origin) in $\R2$ or $\R3$ can be defined geometrically. However, geometry is not as clear in $\R{n}$ for $n > 3$. Thus the angle between two nonzero vectors $x, y \in \R{n}$ is defined to be
\vspace{-.1in}
\[
\arccos \frac{ \ip{x, y} }{ \|x\| \,\|y\| },
\]
where the motivation for this definition comes from Exercise \ref{cost}.
Explain why the Cauchy--Schwarz inequality is needed to show that this definition makes sense.


\exercise
Prove that
\[
\paren[\bigg]{ \sum_{k=1}^n a_k b_k }^2
\le \paren[\bigg]{ \sum_{k=1}^n  k {a_k}^{\2} }  \paren[\bigg]{\sum_{k=1}^n \frac{{b_k}^{\2}}{k} }
\]
for all real numbers $a_1, \dots, a_n$ and $b_1, \dots, b_n$.


\exercise
\begin{SUBtheorem}
\item
Suppose $\fun f {[1, \infty)} {[0, \infty)}$ is continuous. Show that
\[
\paren[\Big]{ \int_1^\infty f\,}^2  \le \int_1^\infty x^2 \paren[\big]{ f(x) }^2 \, dx.
\]
\item
For which continuous functions $\fun f {[1, \infty)} {[0, \infty)}$ is the inequality in (a) an equality with both sides finite?
\end{SUBtheorem}


\exercise
Suppose $v_1, \ldots, v_n$ is a basis of $V$ and $T \in \L(V)$. Prove that if $\la$ is an eigenvalue of $T$, then \label{2HSeigen}
\[
\abs{\la}^2 \le \sum_{\js}^n \sum_{k=1}^n \abs{\M(T)_{j,k}}^2\1,
\]
where $\M(T)_{j,k}$ denotes the entry in row $j$, column $k$ of the matrix of $T$ with respect to the basis $v_1, \ldots, v_n$.


\exercise
Prove that if $u, v \in V\1$, then $\abs[\Big]{\, \norm{u} - \norm{v} \,} \le \norm{u-v}$.\label{6reverse}\index{reverse triangle inequality}
\exercisecomment{The inequality above is called the \emph{reverse triangle inequality}. For the reverse triangle inequality when\/
$V = \C{}$, see Exercise \ref{4reverse} in Chapter \ref{173}.}


\exercise
Suppose $u, v \in V$ are such that
\[
\|u\| = 3, \quad \|u + v\| = 4, \quad \|u - v\| = 6.
\]
What number does $\|v\|$ equal?


\exercise
Show that if $u, v \in V\1$, then
\[
\|u + v\| \, \|u - v\| \le \|u\|^2 + \|v\|^2\1.
\]
\exercisedisplayend


\exercise
Suppose $v_1, \dots, v_m \in V$ are such that $\norm{v_k} \le 1$ for each $k= 1, \ldots, m$. Show that there exist
$a_1, \dots, a_m \in \{1, -1\}$ such that
\[
\norm{a_1 v_1 + \dots + a_m v_m} \le \sqrt{m}.
\]
\exercisedisplayend


\exercise
Prove or give a counterexample: If $\norm{\Cdot}$ is the norm associated with an inner product on $\RR{2}$, then there exists $(x, y) \in \RR{2}$ such that
$
\norm{(x, y)} \ne \max\{|x|, |y|\}$.


\exercise
Suppose $p > 0$. Prove that there is an inner product on $\R{2}$ such that the associated norm is given by
\[
\|(x, y)\| = \paren[\big]{|x|^p + |y|^p}^{1/p}
\]
for all $(x, y) \in \R{2}$ if and only if $p = 2$.


\exercise
Suppose $V$ is a real inner product space. Prove that \label{201}
\[
\ip{u, v} = \frac{\|u + v\|^2 - \|u - v\|^2}{4}
\]
for all $u, v \in V\1$.


\exercise
Suppose $V$ is a complex inner product space. Prove that \label{202}
\[
\ip{u, v} =
\frac{\|u + v\|^2 - \|u - v\|^2 + \|u + iv\|^2 i - \|u -iv\|^2 i}{4}
\]
for all $u, v \in V\1$.


\exercise
A norm on a vector space $U$ is a function
\[
\norm{ \Cdot } \colon U \rightarrow [0, \infty)
\]
such that $\|u\| = 0$ if and only if
$u = 0$, $\|\al u\| = |\al| \|u\|$ for all
$\al \in \F{}$ and all $u \in U$, and $\|u + v\| \le \|u\| + \|v\|$ for all
$u, v \in U$.
Prove that a norm satisfying  the parallelogram equality comes from  an
inner product \big(in other words, show that if $\norm{\Cdot}$ is a norm  on~$U$
satisfying the parallelogram equality, then there is an inner product
$\ip{\cdot ,\cdot }$ on $U$ such that $\|u\| = \ip{u,u}^{1/2}$ for all
$u \in U$\big).


\exercise
Suppose $V\1_1, \dots, V\1_m$ are inner product spaces. Show that the equation
\[
\ip[\big]{(u_1, \dots, u_m), (v_1, \dots, v_m)} =
\ip{u_1, v_1} + \dots + \ip{u_m, v_m}
\]
defines an inner product on $V\1_1 \times \dots \times V\1_m$.
\exercisecomment{In the expression above on the right, for each $k = 1, \ldots, m$, the inner product $\ip{u_k, v_k}$ denotes the inner product on $V\1_k$. Each of the spaces $V\1_1, \dots, V\1_m$ may have a different inner product, even though the same notation is used here.}


\exercise
Suppose $V$ is a real inner product space. For $u, v, w, x\in V\1$, define\index{complexification!of an inner product space}
\[
\ip{u+iv, w + ix}_{\C{}} = \ip{u,w} + \ip{v,x} + (\ip{v,w} - \ip{u,x})i.
\]
\begin{subtheorem}*
\item
Show that $\ip{\cdot, \cdot}_{\C{}}$ makes $V\1_{\C{}}$ into a complex inner product space.
\item
Show that if $u, v \in V\1$, then
\[
\ip{u, v}_{\C{}} = \ip{u, v} \andd {\norm{u + iv}_{\C{}}}^{\2} = \norm{u}^2 + \norm{v}^2\1.
\]
\end{subtheorem}
\exercisecomment{See Exercise \ref{1VC} in Section \ref{1DefVS} for the definition of the complexification\/ $V\1_{\C{}}$.}


\exercise
Suppose $u, v, w \in V\1$. Prove that\label{6wuv}
\[
\norm[\big]{w - \tfrac{1}{2}(u+v)}^2 = \frac{\|w-u\|^2 + \|w-v\|^2}{2} -\frac{\|u-v\|^2}{4}.
\vspace{-.05in}
\]


\exercise
Suppose that $E$ is a subset of $V$ with the property that $u , v \in E$ implies $\frac{1}{2}(u + v) \in E$. Let $w \in V\1$. Show that there is at most one point in $E$ that is closest to $w$. In other words, show that there is at most one $u \in E$ such that
\[
\|w - u\| \le \|w - x\|
\]
for all $x \in E$.


\exercise
Suppose $f, g$ are differentiable functions from $\R{}$ to $\RR{n}$.
\begin{subtheorem}
\item
Show that
\[
\ip[\big]{ f(t), g(t) }' = \ip[\big]{ f' \! (t), g(t)} + \ip[\big]{f(t), g' \! (t) }.
\]
\item
Suppose $c$ is a positive number and $\norm[\big]{f(t)} = c$ for every $t \in \R{}$. Show that $\ip[\big]{ f' \! (t), f(t) } = 0$ for every $t \in \R{}$.
\item
Interpret the result in (b) geometrically in terms of the tangent vector to a curve lying on a sphere in $\R{n}$ centered at the origin.
\end{subtheorem}
\exercisecomment{A function $f \colon \R{} \to \R{n}$ is called differentiable if there exist differentiable functions $f_1, \dots, f_n$ from $\R{}$ to $\R{}$ such that $f(t) =  \bigl(f_1(t), \dots, f_n(t)\bigr)$ for each $t \in \R{}$. Furthermore, for each $t \in \R{}$, the derivative $f' \! (t) \in \R{n}$ is defined by $f' \! (t) = \bigl({f_1}' \! (t), \dots, {f_n}' \! (t)\bigr)$.}


\exercise
Use inner products to prove Apollonius's identity: In a triangle with sides of length $a$, $b$, and $c$, let $d$ be the length of the line segment from the midpoint of the side of length $c$ to the opposite vertex. Then \label{Apollonius}\index{Apollonius's identity}
\[
a^2 + b^2 = \tfrac{1}{2} c^2 + 2d^2\1.
\]
\vspace{-.45in}
\FigureHereNoCaption{Mathematica/Apollonius.pdf}{1}


\exercise
Fix a positive integer $n$. The \emph{Laplacian}\index{Laplacian} $\Delta p$\sindex{Delta@$\Delta$} of a twice differentiable real-valued function $p$ on $\R{n}$ is the function on $\R{n}$ defined by \label{harmonic}
\[
\Delta p = \frac{ \partial^2 p }{\partial {x_1}^{\2}} + \dots + \frac{ \partial^2 p }{\partial {x_n}^{\2}}.
\]
The function $p$ is called \emph{harmonic} if $\Delta p = 0$.\index{harmonic function}

A \emph{polynomial} on $\R{n}$ is a linear combination (with coefficients in $\R{}$) of functions of the form
${x_1}^{\2[m_1]} \dotsm {x_n}^{\2[m_n]}$, where $m_1, \dots, m_n$ are nonnegative integers.

Suppose $q$ is a polynomial on $\RR{n}$. Prove that there exists a harmonic polynomial $p$ on $\R{n}$ such that $p(x) = q(x)$ for every $x \in \R{n}$ with $\|x\| = 1$.
\exercisecomment{The only fact about harmonic functions that you need for this exercise is that if $p$ is a harmonic function on $\R{n}$ and $p(x) = 0$ for all $x \in \R{n}$ with $\|x\|=1$, then $p = 0$.}
\hint{A reasonable guess is that the desired harmonic polynomial $p$ is of the form $q + \paren[\big]{1 - \|x\|^2} r$ for some polynomial $r$. Prove that there is a polynomial $r$ on $\R{n}$ such that $q + \paren[\big]{1 - \|x\|^2} r$ is harmonic by defining an operator $T$ on a suitable vector space by
\[
Tr = \Delta \bigl( (1 - \|x\|^2) r \bigr)
\]
and then showing that $T$ is injective and hence surjective.}


\end{exercises}

\QuoteAuthor{%
In realms of numbers, where the secrets lie,\\
A noble truth emerges from the deep,\\
Cauchy and Schwarz, their wisdom they apply,\\
An inequality for all to keep.\\[-5bp]
\ \\
Two vectors, by this bond, are intertwined,\\
As inner products weave a gilded thread,\\
Their magnitude, by providence, confined,\\
A bound to which their destiny is wed.\\[-5bp]
\ \\
Though shadows fall, and twilight dims the day,\\
This inequality will stand the test,\\
To guide us in our quest, to light the way,\\
And in its truth, our understanding rest.\\[-5bp]
\ \\
So sing, ye muses, of this noble feat,\\
Cauchy--Schwarz, the bound that none can beat.}
{written by ChatGPT with input \textit{Shakespearean sonnet on Cauchy--Schwarz \mbox{inequality}}\index{ChatGPT}}

\section{Orthonormal Bases}  \label{orthobases}

\subsection{Orthonormal Lists and the Gram--Schmidt Procedure}

\DefinitionBox{orthonormal}{\begin{Itemize}
\item
A list of vectors is called \emph{orthonormal}\index{orthonormal!list} if each vector in the list has norm~$1$ and is orthogonal to all the other vectors in the list.
\item
In other words, a list $e_1, \dots, e_m$ of vectors in $V$ is  orthonormal if
\[
\ip{e_j, e_k} = \begin{cases}
1 &\text{if $j = k$,}\\
0 &\text{if $j \ne k$}
\end{cases}
\]
for all $j, k \in \br{1, \ldots, m}$.
\end{Itemize}}

\begin{example}{orthonormal lists}
\begin{subtheorem}*
\item
The standard basis of~$\F{n}$ is an orthonormal list.

\item
$\paren[\Big]{ \frac{1}{\sqrt{3}}, \frac{1}{\sqrt{3}}, \frac{1}{\sqrt{3}} },
\paren[\Big]{ -\frac{1}{\sqrt{2}}, \frac{1}{\sqrt{2}}, 0 } $ is an orthonormal list in $\FF3$.

\item
$\paren[\Big]{ \frac{1}{\sqrt{3}}, \frac{1}{\sqrt{3}}, \frac{1}{\sqrt{3}} },
\paren[\Big]{ -\frac{1}{\sqrt{2}}, \frac{1}{\sqrt{2}}, 0 },
\paren[\Big]{\frac{1}{\sqrt{6}}, \frac{1}{\sqrt{6}}, -\frac{2}{\sqrt{6}} } $ is an orthonormal list in $\FF3$.

\item
Suppose $n$ is a positive integer.  Then, as Exercise \ref{8Fourier} asks you to verify,
\[
\frac{1}{\sqrt{2\pi}}, \frac{\cos x}{\sqrt{\pi}}, \frac{\cos 2x}{\sqrt{\pi}}, \dots,
\frac{\cos nx}{\sqrt{\pi}}, \frac{\sin x}{\sqrt{\pi}}, \frac{\sin 2x}{\sqrt{\pi}}, \dots,
\frac{\sin nx}{\sqrt{\pi}}
\]
is an orthonormal list of vectors in $C[-\pi,\pi]$, the vector space of
continuous real-valued functions on $[-\pi,\pi]$ with inner product
\[
\ip{f, g} = \int_{-\pi}^{\pi} fg.
\]
The orthonormal list above is often used for modeling periodic phenomena, such as tides.

\item
Suppose we make $\P\1_2(\R{})$ into an inner product space using the inner product given by
\[
\ip{p, q} = \int_{-1}^{1} pq
\]
for all $p, q \in \P\1_2(\R{})$. The standard basis $1, x, x^2$ of $\P\1_2(\R{})$ is not an orthonormal list because the vectors in that list do not have norm $1$. Dividing each vector by its norm gives the list $1/\!\sqrt{2}, \sqrt{3/2} x, \sqrt{5/2}x^2\1$, in which each vector has norm $1$, and the second vector is orthogonal to the first and third vectors. However, the first and third vectors are not orthogonal. Thus this is not an orthonormal list. Soon we will see how to construct an orthonormal list from the standard basis
$1, x, x^2$ (see Example \ref{obP2R}).
\end{subtheorem}
\exampleitemend
\end{example}

\pagebreak[3]

Orthonormal lists are particularly easy to work with, as illustrated by the next result.

\ResultBox[4strange]{norm of an orthonormal linear combination}{%
Suppose $e_1, \dots, e_m$ is an orthonormal list of vectors in~$V\1$. Then
\[
\| a_1 e_1 + \dots + a_m e_m \|^2 = | a_1|^2 + \dots + |a_m|^2
\]
for all $a_1, \dots, a_m \in \F{}\3$.
}

\proof
Because each $e_k$ has norm $1$, this follows from repeated applications of the Pythagorean theorem (\ref{184}).
\qed

The result above has the following important corollary.

\ResultBox[46]{orthonormal lists are linearly independent}{Every orthonormal list of vectors is linearly independent.}

\proof
Suppose $e_1, \dots, e_m$ is an orthonormal list of vectors in~$V$ and
$a_1, \dots, a_m \in \F{}$ are such that
\[
a_1 e_1 + \dots + a_m e_m = 0.
\]
Then $| a_1|^2 + \dots + |a_m|^2 = 0$ (by \ref{4strange}), which means that all the
$a_k$'s are $0$. Thus $e_1, \dots, e_m$ is linearly independent.
\qed

%\mar[3]{The next result is named for Friedrich Bessel (1784--1846), who discovered a special case of this inequality.}

Now we come to an important inequality.

\medskip

\ResultBox[Bessel]{Bessel's inequality}
{Suppose $e_1, \dots, e_m$ is an orthonormal list of vectors in~$V\1$. If $v \in V$ then\index{Bessel's inequality}
\[
\abs[\big]{\ip{v, e_1}}^2 + \cdots + \abs[\big]{\ip{v, e_m}}^2 \le \norm{v}^2\1.
\]
}

\proof
Suppose $v \in V\1$. Then
\[
v =\underbrace{\ip{v, e_1} e_1 + \cdots + \ip{v, e_m} e_m }_u +
\underbrace{ v - \ip{v, e_1} e_1 - \cdots - \ip{v, e_m} e_m }_w\!.
\]

Let $u$ and $w$ be defined as in the equation above.
If $k \in \br{1, \ldots, m}$, then $\ip{w, e_k} = \ip{v, e_k} - \ip{v, e_k} \ip{e_k, e_k} = 0$. This implies that $\ip{w, u} = 0$. The Pythagorean theorem now implies that
\begin{align*}
\norm{v}^2 &= \norm{u}^2 + \norm{w}^2 \\[3bp]
&\ge \norm{u}^2 \\[3bp]
&= \abs[\big]{\ip{v, e_1}}^2 + \cdots + \abs[\big]{\ip{v, e_m}}^2\1,
\end{align*}
where the last line comes from \ref{4strange}.
\qed

\pagebreak[2]

The next definition introduces one of the most useful concepts in the study of inner product spaces.

\DefinitionBox{orthonormal basis}{An \emph{orthonormal basis}\index{orthonormal!basis} of $V$ is an orthonormal list of vectors in~$V$ that is also a basis of~$V\1$.}

For example, the standard basis is an orthonormal basis of~$\FF{n}$.

\ResultBox[onbasis]{orthonormal lists of the right length are orthonormal bases}
{Suppose $V$ is finite-dimensional. Then every orthonormal list of vectors in~$V$ of length $\dim V$ is an orthonormal basis of~$V\1$.}

\proof
By \ref{46}, every orthonormal list of vectors in $V$ is linearly independent. Thus every such list of the right length is a basis---see~\ref{75}.
\qed

\begin{example}[4FonBasis]{an orthonormal basis of\/ $\F{4}$}
As mentioned above,  the standard basis is an orthonormal basis of~$\FF{4}$. We now show that
\[
\paren[\Big]{\tfrac{1}{2}, \tfrac{1}{2}, \tfrac{1}{2}, \tfrac{1}{2}},
\paren[\Big]{\tfrac{1}{2}, \tfrac{1}{2}, -\tfrac{1}{2}, -\tfrac{1}{2}},
\paren[\Big]{\tfrac{1}{2}, -\tfrac{1}{2}, -\tfrac{1}{2}, \tfrac{1}{2}},
\paren[\Big]{-\tfrac{1}{2}, \tfrac{1}{2}, -\tfrac{1}{2}, \tfrac{1}{2}}
\]
is also an orthonormal basis of $\FF{4}$.

We have
\[
\norm[\Big]{ \paren[\Big]{\tfrac{1}{2}, \tfrac{1}{2}, \tfrac{1}{2}, \tfrac{1}{2}} }
= \sqrt{ \tfrac{1}{2^2} + \tfrac{1}{2^2} + \tfrac{1}{2^2}+ \tfrac{1}{2^2} }  = 1.
\]
Similarly, the other three vectors in the list above also have norm $1$.

Note that
\[
\ip[\Big]{ \paren[\Big]{\tfrac{1}{2}, \tfrac{1}{2}, \tfrac{1}{2}, \tfrac{1}{2}},
\paren[\Big]{ \tfrac{1}{2}, \tfrac{1}{2}, -\tfrac{1}{2}, -\tfrac{1}{2}} } =
\tfrac{1}{2} \cdot \tfrac{1}{2} + \tfrac{1}{2} \cdot \tfrac{1}{2}
+ \tfrac{1}{2} \cdot \paren[\Big]{-\tfrac{1}{2}} + \tfrac{1}{2} \cdot \paren[\Big]{ -\tfrac{1}{2} } = 0.
\]
Similarly, the inner product of any two distinct vectors in the list above also equals~$0$.

Thus the list above is orthonormal. Because we have an orthonormal list of length four in the four-dimensional vector space $\FF{4}$, this list is an orthonormal basis of $\F{4}$ (by \ref{onbasis}).
\end{example}

In general, given a basis $e_1, \dots, e_n$ of $V$ and a vector $v \in V\1$, we
know that there is some choice of scalars $a_1, \dots, a_n \in \F{}$ such that
\[
v = a_1 e_1 + \dots + a_n e_n.
\]

\noindent Computing the numbers $a_1, \dots, a_n$ that satisfy the equation above can be a long computation for an arbitrary basis of $V\1$. The next result shows, however,
that this is easy for an orthonormal basis---just take
$
a_k = \ip{v, e_k}$.

\mar{%The importance of orthonormal bases stems mainly from the next result.
The formula below for\/ $\norm{v}$ is called \mbox{Parseval's} identity. It was published in 1799 in the context of Fourier \mbox{series}.\index{Parseval's identity}}

Notice how the next result makes each inner product space of dimension $n$ behave like $\FF{n}$, with the role of the \mbox{coordinates} of a vector in $\F{n}$ played by $\ip{v, e_1}, \ldots, \ip{v, e_n}$.

\ResultBox[190]{writing a vector as a linear combination of an orthonormal basis}{%
Suppose\/ $e_1, \dots, e_n$ is an orthonormal basis of\/ $V$ and $u, v \in V\1$.  Then
\begin{subtheorem}*
\item
$v = \ip{v, e_1} e_1 + \dots + \ip{v, e_n} e_n$;
\item
$\|v\|^2 = \abs[\big]{ \ip{v, e_1} }^2  + \dots + \abs[\big]{ \ip{v, e_n} }^2\1$;
\item
$\ip{u, v} = \ip{u, e_1} \overline{ \ip{v, e_1}} + \cdots + \ip{u, e_n} \overline{ \ip{v, e_n}}$.
\end{subtheorem}
}

\proof
Because $e_1, \dots, e_n$ is a basis of~$V\1$, there exist
scalars $a_1, \dots, a_n$ such that
\[
v = a_1 e_1 + \dots + a_n e_n.
\]
Because $e_1, \dots, e_n$ is orthonormal, taking the inner product of both sides of this equation with $e_k$ gives
\mbox{$\ip{v, e_k} = a_k$}.  Thus (a) holds.

Now (b) follows immediately from (a) and~\ref{4strange}.

Take the inner product of $u$ with each side of (a) and then get (c) by using conjugate linearity [\ref{basicip}(d) and \ref{basicip}(e)] in the second slot of the inner product.
\qed

\begin{example}{finding coefficients for a linear combination}
Suppose we want to write the vector $(1,2,4,7) \in \F4$ as a linear combination of the orthonormal basis
\[
\paren[\Big]{\tfrac{1}{2}, \tfrac{1}{2}, \tfrac{1}{2}, \tfrac{1}{2}},
\paren[\Big]{\tfrac{1}{2}, \tfrac{1}{2}, -\tfrac{1}{2}, -\tfrac{1}{2}},
\paren[\Big]{\tfrac{1}{2}, -\tfrac{1}{2}, -\tfrac{1}{2}, \tfrac{1}{2}},
\paren[\Big]{-\tfrac{1}{2}, \tfrac{1}{2}, -\tfrac{1}{2}, \tfrac{1}{2}}
\]
of $\F4$ from Example \ref{4FonBasis}. Instead of solving a system of four linear equations in four unknowns, as typically would be required if we were working with a nonorthonormal basis, we simply evaluate four inner products and use \ref{190}(a), getting that
$(1,2,4,7)$ equals
\[
7 \paren[\Big]{\tfrac{1}{2}, \tfrac{1}{2}, \tfrac{1}{2}, \tfrac{1}{2}}
-4\paren[\Big]{\tfrac{1}{2}, \tfrac{1}{2}, -\tfrac{1}{2}, -\tfrac{1}{2}}
+\paren[\Big]{\tfrac{1}{2}, -\tfrac{1}{2}, -\tfrac{1}{2}, \tfrac{1}{2}}
+2\paren[\Big]{-\tfrac{1}{2}, \tfrac{1}{2}, -\tfrac{1}{2}, \tfrac{1}{2}}.
\]
\exampledisplayend
\end{example}

Now that we understand the usefulness of orthonormal bases, how do we
go about finding them?  For example, does $\P\1_m(\R{})$ with inner
product as in \ref{exampleip}(c) have an
orthonormal basis? The next result will lead to
answers to these questions.

\mar[5]{J{\o}rgen Gram \uppar{1850--1916}\index{Gram, J{\o}rgen} and \mbox{Erhard} Schmidt \uppar{1876--1959}\index{Schmidt, Erhard} popularized this algorithm that constructs orthonormal lists.}

The algorithm used in the next proof is
called the \emph{Gram--Schmidt procedure}.\index{Gram--Schmidt procedure}
It gives a method for turning a linearly independent list into an orthonormal
list with the same span as the original list.

\ResultBox[186]{Gram--Schmidt procedure}{%
Suppose $v_1, \dots, v_m$ is a linearly independent list of vectors in\/~$V\1$.
Let \mbox{$f_1 = v_1$}. For $k = 2, \dots, m$, define $f_k$ inductively by
\[
f_k = v_k - \frac{\ip{v_k, f_1}}{\norm{f_1}^2} f_1 - \dots - \frac{\ip{v_k, f_{k-1} }}{\norm{f_{k-1}}^2} f_{k-1}.%
\]
For each $k = 1, \ldots, m$, let \raisebox{0bp}[9bp][0bp]{$e_k = \dfrac{f_k}{\norm{f_k}}$}. Then $e_1, \dots, e_m$ is an orthonormal list of vectors in\/~$V$ such that
\[
\Span (v_1, \dots, v_k) = \Span (e_1, \dots, e_k)
\]
for each $k = 1, \dots, m$.
}

\proof
We will show by induction on $k$ that the desired conclusion holds. To get started with $k = 1$, note that because $e_1 = f_1/\norm{f_1}$, we have
$\norm{e_1} = 1$; also, $\Span(v_1) = \Span(e_1)$ because $e_1$ is a nonzero multiple of $v_1$.

Suppose $1 < k \le m$ and the list $e_1, \dots, e_{k-1}$ generated by \ref{186} is an orthonormal list such that
\begin{equation} \label{83}
\Span (v_1, \dots, v_{k-1}) = \Span (e_1, \dots, e_{k-1}).
\end{equation}%
Because $v_1, \dots, v_m$ is linearly independent, we have $v_k \notin \Span (v_1, \dots, v_{k-1})$. Thus
$v_k \notin \Span (e_1, \dots, e_{k-1}) = \Span (f_1, \dots, f_{k-1})$, which implies that $f_k \ne 0$.
Hence  we are not dividing by $0$ in the definition of $e_k$ given in \ref{186}.  Dividing a vector by its norm produces a new vector with norm~$1$; thus $\|e_k\| = 1$.

Let $j \in \br{1, \ldots, k-1}$.  Then
\begin{align*}
\ip{e_k, e_j} &= \frac{1}{\norm{f_k} \, \norm{f_j}} \ip{f_k, f_j}\\[6bp]
&= \frac{1}{\norm{f_k}  \, \norm{f_j}} \ip[\Bigg]{  v_k - \frac{\ip{v_k, f_1}}{\norm{f_1}^2} f_1 - \dots - \frac{\ip{v_k, f_{k-1} }}{\norm{f_{k-1}}^2} f_{k-1}, f_j } \\[6bp]
&= \frac{1}{\norm{f_k}  \, \norm{f_j}} \paren[\big]{\ip{v_k, f_j} - \ip{v_k, f_j}} \\[6bp]
&= 0.
\end{align*}
Thus $e_1, \dots, e_k$ is an orthonormal list.

From the definition of $e_k$ given in \ref{186}, we see that $v_k \in \Span (e_1, \dots, e_k)$.  Combining this
information with~\ref{83} shows that
\[
\Span (v_1, \dots, v_k) \subset \Span (e_1, \dots, e_k).
\]
Both lists above are linearly independent (the $v$'s by hypothesis, and the $e$'s by orthonormality and~\ref{46}).  Thus both subspaces above have dimension~$k$, and
hence they are equal, completing the induction step and thus completing the proof.
\qed

\pagebreak[3]

\begin{example}[obP2R]{an orthonormal basis of\/ $\P\1_2(\R{})$}
Suppose we make $\P\1_2(\R{})$ into an inner product space using the inner product given by
\[
\ip{p, q} = \int_{-1}^{1} pq
\]
for all $p, q \in \P\1_2(\R{})$. We know that $1, x, x^2$ is a basis of $\P\1_2(\R{})$, but it is not an orthonormal basis. We will find an orthonormal basis of $\P\1_2(\R{})$ by applying the Gram--Schmidt procedure with $v_1 = 1$, $v_2 = x$, and $v_3 = x^2\1$.

To get started, take $f_1 = v_1 = 1$. Thus $\norm{f_1}^2 = \int_{-1}^1 1 = 2$. Hence the formula in \ref{186} tells us that
\[
f_2 = v_2 - \frac{\ip{v_2, f_1}}{\norm{f_1}^2} f_1 = x - \frac{\ip{x, 1}}{\norm{f_1}^2} = x,
\]
where the last equality holds because $\ip{x, 1} = \int_{-1}^1 t \, dt = 0$.

The formula above for $f_2$ implies that $\norm{f_2}^2 = \int_{-1}^1 t^2 \, dt = \frac{2}{3}$. Now the formula in \ref{186} tells us that
\[
f_3 = v_3 - \frac{\ip{v_3, f_1}}{\norm{f_1}^2} f_1 - \frac{\ip{v_3, f_2}}{\norm{f_2}^2} f_2
= x^2 - \tfrac{1}{2} \ip[\big]{x^2\1, 1} - \tfrac{3}{2}\ip[\big]{x^2\1, x} x
= x^2 - \tfrac{1}{3}.
\]

The formula above for $f_3$ implies that
\[
\norm{f_3}^2 = \int_{-1}^1 \paren[\Big]{t^2 - \tfrac{1}{3} }^2 \, dt =  \int_{-1}^{1} \paren[\Big]{ t^4 -\tfrac{2}{3} t^2 + \tfrac{1}{9} } \, dt = \tfrac{8}{45}.
\]

Now dividing each of $f_1, f_2, f_3$ by its norm gives us the orthonormal list
\[
\sqrt{\tfrac{1}{2}}, \sqrt{\tfrac{3}{2}} x,
\sqrt{\tfrac{45}{8}} \paren[\Big]{ x^2 - \tfrac{1}{3} }.
\]
The orthonormal list above has length three, which is the dimension of $\P\1_2(\R{})$. Hence this orthonormal list is an orthonormal basis of~$\P\1_2(\R{})$ [by \ref{onbasis}].
\end{example}

Now we can answer the question about the existence of orthonormal bases.

\ResultBox{existence of orthonormal basis}{
Every finite-dimensional inner product space has an orthonormal basis.}

\proof
Suppose $V$ is finite-dimensional. Choose a basis of~$V\1$.  Apply the Gram--Schmidt procedure (\ref{186}) to
it, producing an orthonormal list of length $\dim V\1$. By \ref{onbasis}, this orthonormal list is an orthonormal basis of~$V\1$.
\qed

Sometimes we need to know not only that an orthonormal
basis exists, but also that every orthonormal list can be extended to an
orthonormal basis.  In the next corollary, the Gram--Schmidt procedure shows
that such an extension is always possible.

\ResultBox[49]{every orthonormal list extends to an orthonormal basis}{%
Suppose $V$ is finite-dimensional. Then every orthonormal list of vectors in $V$ can be extended to an orthonormal basis of $V\1$.}

\proof
Suppose $e_1, \dots, e_m$ is an orthonormal list of vectors in $V\1$.  Then
$e_1, \dots, e_m$ is linearly independent (by~\ref{46}). Hence this list can be
extended to a basis $e_1, \dots, e_m, v_1, \dots, v_n$ of~$V$ (see~\ref{9}). Now
apply  the Gram--Schmidt procedure (\ref{186}) to
$e_1, \dots, e_m, v_1, \dots, v_n$, producing an  orthonormal list
\[
e_1, \dots, e_m, f_1, \dots, f_n;
\]
here the formula given by the Gram--Schmidt procedure leaves the first $m$ vectors unchanged because they are already orthonormal.  The list above is an  orthonormal basis of $V$ by \ref{onbasis}.
\qed

Recall that a matrix is called upper triangular if it looks like this:
\[
\mat{ccc}{ * & & * \\ & \ddots & \\ 0 & & * },
\]
where the $0$ in the matrix above indicates that all entries below the diagonal equal~$0$, and asterisks are used to denote entries on and above the diagonal.

In the last chapter, we gave a necessary and sufficient condition for an operator to have an upper-triangular matrix with respect to some basis (see \ref{5necsuf}). Now that we are dealing with inner product spaces, we would like to know whether there
exists an \textit{orthonormal} basis with respect to which we have an upper-triangular matrix. The next result shows that the condition for an operator to have an upper-triangular matrix with respect to some orthonormal basis is the same as the condition to have an upper-triangular matrix with respect to an arbitrary basis.

\ResultBox[188]{upper-triangular matrix with respect to some orthonormal basis}
{Suppose $V$ is finite-dimensional and $T \in \L(V)$. Then $T$ has an upper-triangular matrix with respect to some orthonormal basis of $V$ if and only if the minimal polynomial of $T$ equals $(z-\la_1) \cdots (z-\la_m)$ for some
$\la_1, \ldots, \la_m \in \F{}\3$.\index{minimal polynomial!and upper-triangular matrices}}

\proof
Suppose $T$ has an upper-triangular matrix with respect to some
basis $v_1, \dots, v_n$ of~$V\1$.  Thus $\Span(v_1, \dots, v_k)$ is
invariant under $T$ for each $k = 1, \dots, n$ (see~\ref{a225}).

Apply the Gram--Schmidt procedure to $v_1, \dots, v_n$, producing an
orthonormal basis $e_1, \dots, e_n$ of~$V\1$.  Because
\[
\Span(e_1, \dots, e_k) = \Span(v_1, \dots, v_k)
\]
for each $k$ (see~\ref{186}), we conclude that
$\Span(e_1, \dots, e_k)$ is invariant under~$T$ for each $k = 1, \dots, n$.  Thus,
by~\ref{a225}, $T$ has an upper-triangular matrix with respect to the orthonormal basis
$e_1, \dots, e_n$. Now use \ref{5necsuf} to complete the proof.
\qed

\mar{Issai Schur\index{Schur, Issai} \uppar{1875--1941} published a proof of the next result in 1909.}

For complex vector spaces, the next result is an important application of the result above. See Exercise \ref{3Schur} for a version of Schur's theorem that applies simultaneously to more than one operator.

\ResultBox[214]{Schur's theorem\index{Schur's theorem}}
{Every operator on a finite-dimensional complex inner product space has an upper-triangular matrix with respect to some orthonormal basis.}

\proof
The desired result follows from the second version of the fundamental theorem of algebra (\ref{4FTA2}) and \ref{188}.
\qed

\subsection{Linear Functionals on Inner Product Spaces}

Because linear maps into the scalar field $\F{}$ play a special role, we defined a special name for them and their vector space in Section \ref{duals}. Those definitions are repeated below in case you skipped Section \ref{duals}.

\DefinitionBox{linear functional, dual space, $V'$\sindex{V'@$V'\1$}}
{\begin{Itemize}
\item
A \emph{linear functional}\index{linear functional} on $V$ is a linear map from $V$ to $\F{}\3$.
\item
The \emph{dual space}\index{dual!of a vector space} of $V\1$, denoted by $V'\!$, is the vector space of all linear \mbox{functionals} on $V\1$. In other words, $V' = \L(V\1, \F{})$.
\end{Itemize}
}

\begin{example}[359]{linear functional on\/ $\F{3}$}
The function
$\phi \colon \F{3} \rightarrow \F{}$ defined by
\[
\phi(z_1, z_2, z_3) = 2 z_1 - 5 z_2 + z_3
\]
is a linear functional on $\FF{3}$. We could write this linear functional in the form
\[
\phi(z) = \ip{z, w}
\]
for every $z \in \FF3$, where $w = (2, -5, 1)$.
\end{example}

\begin{example}[360]{linear functional on\/ $\P\1_5(\R{})$}
The function $\phi \colon \P\1_5(\R{}) \rightarrow \R{}$ defined by
\[
\phi(p) = \int_{-1}^1 p(t) \bigl(\cos(\pi t)\bigr)\,dt
\]
is a linear functional on $\P\1_5(\R{})$.
\end{example}

\mar{The next result is named in honor of Frigyes Riesz \uppar{1880--1956}, who proved several theorems early in the twentieth century that look very much like the result below.\index{Riesz, Frigyes}}

If $v \in V\1$, then the map that sends $u$ to~$\ip{u,v}$ is a linear functional
on~$V\1$. The next result states  that every linear functional on $V$ is of this form. For example, we can take $v = (2, -5, 1)$ in Example \ref{359}.

Suppose we make the vector space $\P\1_5(\R{})$ into an inner product space by defining $\ip{p, q} = \int_{-1}^{1} pq$. Let $\phi$ be as in Example \ref{360}. It is not obvious that there exists $q \in \P\1_5(\R{})$ such that
\[
\int_{-1}^1 p(t) \bigl(\cos(\pi t)\bigr)\,dt = \ip{p, q}
\]
for every $p \in \P\1_5(\R{})$ \sbr{we cannot take
$q(t) = \cos(\pi t)$ because that choice of $q$ is not an element of~$\P\1_5(\R{})$}. The next result tells us the somewhat surprising result that there indeed exists a polynomial $q \in \P\1_5(\R{})$ such that the equation above holds for all $p \in \P\1_5(\R{})$.

\ResultBox[50]{Riesz representation theorem\index{Riesz representation theorem}}{%
Suppose $V$ is finite-dimensional and $\phi$ is a linear functional on $V\1$.  Then there is a unique vector
$v \in V$ such that
\[
\phi(u) = \ip{u,v}
\]
for every $u \in V\1$.}

\proof
First we show that there exists a vector $v \in V$ such that $\phi(u) = \ip{u, v}$ for every $u \in V\1$.  Let $e_1, \dots, e_n$ be an orthonormal basis of $V\1$.  Then
\begin{align*}
\phi(u) &=
\phi\paren[\big]{\ip{u, e_1} e_1 + \dots + \ip{u, e_n} e_n} \\[3bp]
&= \ip{u, e_1} \phi(e_1) + \dots + \ip{u, e_n} \phi(e_n)  \\[3bp]
&= \ip[\Big]{u, \overline{\phi(e_1)} e_1  + \dots + \overline{\phi(e_n)} e_n}
\end{align*}
for every $u \in V\1$, where the first equality comes from~\ref{190}(a).  Thus
setting
\begin{equation} \label{GSformula}
v = \overline{\phi(e_1)} e_1  + \dots +
\overline{\phi(e_n)} e_n,
\end{equation}%
we have $\phi(u) = \ip{u, v}$ for every $u \in V\1$, as desired.

Now we prove that only one vector $v \in V$ has the desired behavior.  Suppose
$v_1, v_2 \in V$ are such that
\[
\phi(u) = \ip{u, v_1} = \ip{u, v_2}
\]
for every $u \in V\1$.  Then
\[
0 = \ip{u, v_1} - \ip{u, v_2} = \ip{u, v_1 - v_2}
\]
for every $u \in V\1$.  Taking $u = v_1 - v_2$ shows that $v_1 - v_2 =0$. Thus $v_1 = v_2$, completing the proof of the uniqueness part of the result.
\qed

\begin{example}{computation illustrating Riesz representation theorem}
Suppose we want to find a polynomial $q \in \P\1_2(\R{})$ such that
\begin{equation} \label{2rrt}
\int_{-1}^1 p(t) \bigl(\cos(\pi t)\bigr)\,dt
= \int_{-1}^1 pq
\end{equation}%
for every polynomial $p \in \P\1_2(\R{})$. To do this, we make $\P\1_2(\R{})$ into an inner product space by defining $\ip{p, q}$ to be the right side of the equation above for \mbox{$p, q \in \P\1_2(\R{})$}. Note that the left side of the equation above does not equal the inner product in $\P\1_2(\R{})$ of $p$ and the function $t \mapsto \cos(\pi t)$ because this last function is not a polynomial.

Define a linear functional $\phi$ on $\P\1_2(\R{})$ by letting
\[
\phi(p) = \int_{-1}^1 p(t) \bigl(\cos(\pi t)\bigr)\,dt
\]
for each $p \in \P\1_2(\R{})$. Now use the orthonormal basis from Example \ref{obP2R} and apply formula \ref{GSformula} from the proof of the Riesz representation theorem to see that if $p \in \P\1_2(\R{})$, then $\phi(p) = \ip{p, q}$, where
\begin{align*}
q(x) &= \paren[\bigg]{ \int_{-1}^1 \sqrt{\tfrac{1}{2}} \cos(\pi t) \, dt } \sqrt{\tfrac{1}{2}} +
 \paren[\bigg]{ \int_{-1}^1 \sqrt{\tfrac{3}{2}} \, t \cos(\pi t) \, dt } \sqrt{\tfrac{3}{2}}x \\[6bp]
 & \quad +
  \paren[\bigg]{ \int_{-1}^1 \sqrt{\tfrac{45}{8}} \paren[\Big]{ t^2 - \tfrac{1}{3} } \cos(\pi t) \, dt } \sqrt{\tfrac{45}{8}}\paren[\Big]{ x^2 - \tfrac{1}{3} }.
\end{align*}
A bit of calculus applied to the equation above shows that
\[
q(x) = \tfrac{15}{2 \pi^2} \paren[\big]{1 - 3 x^2}.
\]

The same procedure shows that if we want to find $q \in \P\1_5(\R{})$ such that \ref{2rrt} holds for all $p \in \P\1_5(\R{})$, then we should take
\[
q(x) = \tfrac{105}{8 \pi^4} \paren[\bigg]{ \paren[\big]{27 - 2 \pi^2} + \paren[\big]{24 \pi^2 - 270} x^2 + \paren[\big]{315 - 30 \pi^2} x^4}.
\]
\exampledisplayend\vspace{4bp}
\end{example}

Suppose $V$ is finite-dimensional and $\phi$ a linear functional on $V\1$. Then \ref{GSformula} gives a formula for the vector $v$ that satisfies
\[
\phi(u) = \ip{u, v}
\]
for all $u \in V\1$. Specifically, we have
 \[
 v = \overline{\phi(e_1)} e_1  + \dots + \overline{\phi(e_n)} e_n.
\]
The right side of the equation above seems to depend on the orthonormal basis $e_1, \dots, e_n$ as well as on $\phi$. However, \ref{50} tells us that $v$ is uniquely determined by $\phi$. Thus the right side of the equation above is the same regardless of which orthonormal basis $e_1, \dots, e_n$ of $V$ is chosen.

For two additional different proofs of the Riesz representation theorem, see \ref{3RRT} and also Exercise~\ref{VVpiso} in Section~\ref{OrthoMin}.

\begin{exercises}

\exercise
Suppose $e_1, \dots, e_m$ is a list of vectors in~$V$ such that
\[
\| a_1 e_1 + \dots + a_m e_m \|^2 = | a_1|^2 + \dots + |a_m|^2
\]
for all $a_1, \dots, a_m \in \F{}\3$. Show that $e_1, \dots, e_m$ is an orthonormal list.
\exercisecomment{This exercise provides a converse to \ref{4strange}.}


\exercise
\begin{SUBtheorem}
\item
Suppose $\theta \in \R{}$. Show that both
\[
(\cos \theta, \sin \theta), (-\sin \theta, \cos \theta) \andd (\cos \theta, \sin \theta), (\sin \theta, -\cos \theta)
\]
are orthonormal bases of $\RR2$.
\item
Show that each orthonormal basis of $\R2$ is of the form given by one of the two possibilities in (a).
\end{SUBtheorem}


\exercise
Suppose $e_1, \dots, e_m$ is an orthonormal list in~$V$ and $v \in V\1$.  Prove~that
\[
\|v\|^2 = \abs[\big]{ \ip{v, e_1} }^2 + \dots + \abs[\big]{ \ip{v, e_m} }^2 \iff v \in \Span(e_1, \dots, e_m).
\]
\exercisedisplayend


\exercise
Suppose $n$ is a positive integer.  Prove that \label{8Fourier}
\[
\frac{1}{\sqrt{2\pi}}, \frac{\cos x}{\sqrt{\pi}}, \frac{\cos 2x}{\sqrt{\pi}}, \dots,
\frac{\cos nx}{\sqrt{\pi}}, \frac{\sin x}{\sqrt{\pi}}, \frac{\sin 2x}{\sqrt{\pi}}, \dots,
\frac{\sin nx}{\sqrt{\pi}}
\]
is an orthonormal list of vectors in $C[-\pi,\pi]$, the vector space of
continuous real-valued functions on $[-\pi,\pi]$ with inner product
\[
\ip{f, g} = \int_{-\pi}^{\pi} fg.
\]
\hint{The following formulas should help.\normalsize
\begin{align*}
(\sin x)(\cos y) &= \frac{\sin(x-y) + \sin(x+y)}{2}\\[4bp]
(\sin x)(\sin y) &= \frac{\cos(x-y) - \cos(x+y)}{2}\\[4bp]
(\cos x)(\cos y) &= \frac{\cos(x-y) + \cos(x+y)}{2}
\end{align*}}


\exercise
Suppose $\fun f {[-\pi, \pi]} {\R{}}$ is continuous. For each nonnegative integer $k$, define
\[
a_k = \tfrac{1}{\sqrt{\pi}} \int_{-\pi}^{\pi} f(x) \cos(kx) \,dx \andd b_k = \tfrac{1}{\sqrt{\pi}}  \int_{-\pi}^{\pi} f(x) \sin(kx) \,dx.
\]
Prove that
\[
\frac{{a_0}^{\2}}{2} + \sum_{k=1}^\infty \paren[\big]{ {a_k}^{\2} + {b_k}^{\2} } \le \int_{-\pi}^{\pi} f^2\1.
\]
\exercisecomment{The inequality above is actually an equality for all continuous functions $\fun f {[-\pi, \pi]} {\R{}}$. However, proving that this inequality is an equality involves Fourier series techniques beyond the scope of this book.}


\exercise
Suppose $e_1, \dots, e_n$ is an orthonormal basis of $V\1$.
\begin{subtheorem}
\item
Prove that if $v_1, \dots, v_n$ are vectors in $V$ such that
\[
\|e_k - v_k\| < \frac{1}{\sqrt{n}}
\]
for each $k$, then $v_1, \dots, v_n$ is a basis of $V\1$.

\item
Show that there exist $v_1, \ldots, v_n \in V$ such that
\[
\|e_k - v_k\| \le \frac{1}{\sqrt{n}}
\]
for each $k$, but $v_1, \ldots, v_n$ is not linearly independent.
\end{subtheorem}
\exercisecomment{This exercise states in \uppar{a} that an appropriately small perturbation of an orthonormal basis is a basis. Then \uppar{b} shows that the number $1/\!\sqrt{n}$ on the right side of the inequality in \uppar{a} cannot be higher.}


\exercise
Suppose $T \in \L\paren[\big]{\R{3}}$ has an upper-triangular matrix with respect to the basis $(1, 0, 0)$, (1, 1, 1), $(1, 1, 2)$. Find an orthonormal basis of $\R{3}$ with respect to which $T$ has an upper-triangular matrix.


\exercise
Make $\P\1_2(\R{})$ into an inner product space by defining \label{375}
$
\ip{p, q} = \int_0^1 pq
$
for all $p, q \in \P\1_2(\R{})$.
\begin{subtheorem}
\item
Apply the Gram--Schmidt procedure to the basis $1, x, x^2$ to produce an
orthonormal basis of $\P\1_2(\R{})$.
\item
The differentiation operator\index{differentiation linear map} \big(the operator that
takes $p$ to $p'$\big) on $\P\1_2(\R{})$ has an upper-triangular matrix with respect to
the basis $1, x, x^2\1$, which is not an orthonormal basis. Find the matrix of the differentiation operator on $\P\1_2(\R{})$ with respect to the orthonormal basis produced in (a) and verify that this matrix is upper triangular, as expected from the proof of \ref{188}.
\end{subtheorem}
\exercisesubtheoremend


\exercise
Suppose $e_1, \dots, e_m$ is the result of applying the Gram--Schmidt procedure to a linearly independent list $v_1, \dots, v_m$ in $V\1$. Prove that $\ip{v_k, e_k} > 0$ for each $k = 1, \dots, m$.


\exercise
Suppose $v_1, \dots, v_m$ is a linearly independent list in $V\1$. Explain why the orthonormal list produced by the formulas of the Gram--Schmidt procedure (\ref{186}) is the only orthonormal list $e_1, \ldots, e_m$ in $V$ such that $\ip{v_k, e_k} > 0$ and
$\Span(v_1, \ldots, v_k) = \Span(e_1, \ldots, e_k)$ for each $k = 1, \ldots, m$. \label{8GSunique}
\exercisecomment{The result in this exercise is used in the proof of \ref{7QR}.}


\exercise
Find a polynomial $q \in \P\1_2(\R{})$ such that
$
p\paren[\Big]{\tfrac{1}{2}} = \int_0^1 pq$
for every $p \in \P\1_2(\R{})$.\label{6int}


\exercise
Find a polynomial $q \in \P\1_2(\R{})$ such that
\[
\int_0^1 p(x) \cos(\pi x) \,dx = \int_0^1 pq
\]
for every $p \in \P\1_2(\R{})$.


\exercise
Show that a list $v_1, \ldots, v_m$ of vectors in $V$ is linearly dependent if and only if the Gram--Schmidt formula in \ref{186} produces $f_k= 0$ for some $k \in \br{1, \ldots, m}$.\label{GSlindep}%
\exercisecomment{This exercise gives an alternative to Gaussian elimination techniques for determining whether a list of vectors in an inner product space is linearly dependent.}


%\begin{comment}
\exercise
Suppose $V$ is a real inner product space and $v_1, \dots, v_m$ is a
linearly independent list of vectors in~$V\1$.  Prove that there exist exactly
$2^m$ orthonormal lists $e_1, \dots, e_m$ of vectors
in~$V$ such  that
\[
\Span (v_1, \dots, v_k) = \Span (e_1, \dots, e_k)
\]
for all $k \in \{1, \dots, m\}$.

%\end{comment}

\exercise
Suppose $\ip{ \cdot, \cdot}_1$ and $\ip{ \cdot, \cdot}_2$ are inner products on $V$ such that $\ip{u,v}_1 = 0$ if and only if $\ip{u,v}_2= 0$. Prove that there is a positive number $c$ such that $\ip{u, v}_1 = c \ip{u, v}_2$ for every $u, v \in V\1$.
\exercisecomment{This exercise shows that if two inner products have the same pairs of \mbox{orthogonal} vectors, then each of the inner products is a scalar multiple of the other inner product.}


\exercise
Suppose $V$ is finite-dimensional. Suppose $\ip{ \cdot, \cdot}_1$, $\ip{ \cdot, \cdot}_2$ are inner products on $V$ with corresponding norms $\norm{ \Cdot }_1$ and $\norm{\Cdot}_2$. Prove that there exists a positive number $c$ such that
$
\| v \|_1 \le c \| v \|_2
$
for every $v \in V\1$.


\exercise
Suppose $\F{} = \C{}$ and $V$ is finite-dimensional. Prove that if $T$ is an operator on $V$ such that $1$ is the only eigenvalue of $T$ and $\norm{Tv} \le \norm{v}$ for all $v \in V\1$, then $T$ is the identity operator.


\exercise
Suppose $u_1, \dots, u_m$ is a linearly independent list in $V\1$. Show that there exists $v \in V$ such that $\ip{u_k, v} =1$
for all $k \in \{1, \dots, m\}$.


\exercise
Suppose $v_1, \ldots, v_n$ is a basis of $V\1$. Prove that there exists a basis $u_1, \ldots, u_n$ of $V$ such that
\[
\ip{v_j, u_k} = \begin{cases}
0 & \text{if } j \ne k, \\
1 & \text{if } j=k.
\end{cases}
\]
\vspace{-.19in}


\exercise
Suppose $\F{}=\C{}$, $V$ is finite-dimensional, and $\E \subset \L(V)$ is such that \[ST = TS\] for all $S, T \in \E$. Prove that there is an orthonormal basis of $V$ with respect to which every element of $\E$ has an upper-triangular matrix.%
 \label{3Schur}\index{commuting operators}
\exercisecomment{This exercise strengthens Exercise \ref{4ManyOps}\uppar{b} in Section \ref{2diagmat} \uppar{in the context of inner product spaces} by asserting that the basis in that exercise can be chosen to be orthonormal.}


\exercise
Suppose $\F{} = \C{}$, $V$ is finite-dimensional, $T \in \L(V)$, and all eigenvalues of $T$ have absolute value less than $1$. Let $\epsilon > 0$. Prove that there exists a positive integer $m$ such that
$\norm[\big]{T^m v} \le \epsilon \|v\|$ for every $v \in V\1$.


\exercise
Suppose $C[-1, 1]$ is the vector space of continuous real-valued functions on the interval $[-1 ,1]$ with inner product given by \index{Riesz representation theorem}
\[
\ip{f, g} = \int_{-1}^1 fg
\]
for all $f, g \in C[-1, 1]$. Let $\phi$ be the linear functional on $C[-1, 1]$ defined by $\phi(f) = f(0)$. Show that there does not exist $g \in C[-1, 1]$ such that
\[
\phi(f) = \ip{f, g}
\]
for every $f \in C[-1, 1]$.
\exercisecomment{This exercise shows that the Riesz representation theorem \uppar{\ref{50}} does not hold on infinite-dimensional vector spaces without additional hypotheses on $V$ and $\phi$.}


\exercise
For all $u, v \in V\1$, define $d(u, v) = \|u - v\|$.\label{4metric}
\begin{subtheorem}
\item
Show that $d$ is a metric on $V\1$.
\item
Show that if $V$ is finite-dimensional, then $d$ is a complete metric on $V$ (meaning that every Cauchy sequence converges).
\item
Show that every finite-dimensional subspace of $V$ is a closed subset of~$V$ (with respect to the metric $d$\,).
\end{subtheorem}
\exercisecomment{This exercise requires familiarity with metric spaces.}


\end{exercises}

\vfill

\SpecialBox{orthogonality at the Supreme Court}%
{Law professor Richard Friedman presenting a case before the U.S. Supreme Court in 2010:\index{Supreme Court}
\begin{tabbing}
\textit{Mr.~Friedman}: \=I think that issue is entirely orthogonal to the issue here\\
\>because the Commonwealth is acknowledging---\\[1bp]
\textit{Chief Justice Roberts}: I'm sorry. Entirely what?\\[1bp]
\textit{Mr.~Friedman}: Orthogonal. Right angle. Unrelated. Irrelevant.\\[1bp]
\textit{Chief Justice Roberts}: Oh.\\[1bp]
\textit{Justice Scalia}: What was that adjective? I liked that.\\[1bp]
\textit{Mr.~Friedman}: Orthogonal.\\[1bp]
\textit{Chief Justice Roberts}: Orthogonal.\\[1bp]
\textit{Mr.~Friedman}: Right, right.\\[1bp]
\textit{Justice Scalia}: Orthogonal, ooh.
(Laughter.)\\[1bp]
\textit{Justice Kennedy}: I knew this case presented us a problem.
(Laughter.)
\end{tabbing}}

\section{Orthogonal Complements and Minimization Problems} \label{OrthoMin}

\subsection{Orthogonal Complements}

\DefinitionBox{orthogonal complement, $U^\perp$}{If $U$ is a subset of $V\1$, then the \emph{orthogonal
complement}\index{orthogonal!complement} of $U$,
denoted by $U^\perp\1$,\sindex{U@$U^\perp$} is
the  set of all vectors in $V$ that are orthogonal to every vector in $U$:
\[
U^\perp = \br[\big]{v \in V: \ip{u,v} = 0 \mbox{\ for every\ } u \in U}.
\]}

The orthogonal complement $U^\perp$ depends on $V$ as well as on $U$. However, the inner product space $V$ should always be clear from the context and thus it can be omitted from the notation.

\begin{example}{orthogonal complements}
\begin{Itemize}
\item
If $V = \R3$ and $U$ is the subset of $V$ consisting of the single point $(2,3,5)$, then $U^\perp$ is the plane
$\br[\big]{ (x,y,z) \in \R3: 2x+3y+5z = 0}$.
\item
If $V = \R3$ and $U$ is the plane $\br[\big]{ (x,y,z) \in \R3: 2x+3y+5z = 0}$, then $U^\perp$ is the line
$\br{ (2t, 3t, 5t): t \in \R{} }$.
\item
More generally, if $U$ is a plane in $\R3$ containing the origin, then $U^\perp$ is the line containing the origin that is perpendicular to $U$.
\item
If $U$ is a line in $\R3$ containing the origin, then $U^\perp$ is the plane containing the origin that is perpendicular to $U$.
\item
If $V = \F5$ and $U = \br[\big]{ (a,b, 0, 0, 0) \in \F5: a,b \in \F{} }$, then
\[
U^\perp = \br[\big]{(0, 0, x,y,z) \in \F5: x,y,z \in \F{}}.
\]
\item
If $e_1, \ldots, e_m, f_1, \ldots, f_n$ is an orthonormal basis of $V\1$, then
\[
\paren[\big]{\Span( e_1, \ldots, e_m)}^\perp = \Span( f_1, \ldots, f_n).
\]
\end{Itemize}
\exampleitemend
\end{example}

We begin with some straightforward consequences of the definition.

\ResultBox[oc]{properties of orthogonal complement}
{\begin{subtheorem}
\item
If $U$ is a subset of $V\1$, then $U^\perp$ is a subspace of $V\1$.
\item
$\{0\}^\perp = V\1$.
\item
$V^\perp = \{0\}$.
\item
If $U$ is a subset of $V\1$, then $U \cap U^\perp \subset \{0\}$.
\item
If $G$ and $H$ are subsets of $V$ and $G \subset H$, then $H^\perp \subset G^\perp\1$.
\end{subtheorem}}

\pagebreak[3]

\proof
\begin{subtheorem}*
\item
Suppose $U$ is a subset of $V\1$. Then $\ip{u,0} = 0$ for every $u \in U$; thus $0 \in U^\perp\1$.

Suppose $v, w \in U^\perp\1$. If $u \in U$, then
\[
\ip{u,v+w} = \ip{u,v} + \ip{u,w} = 0 + 0 = 0.
\]
Thus $v + w \in U^\perp\1$, which shows that $U^\perp$ is closed under addition.

Similarly, suppose $\la \in \F{}$ and $v \in U^\perp\1$. If $u \in U$, then
\[
\ip{u, \la v} = \overline{\la} \ip{u,v} = \overline{\la} \cdot 0 = 0.
\]
Thus $\la v \in U^\perp\1$, which shows that $U^\perp$ is closed under scalar multiplication. Thus $U^\perp$ is a subspace of $V\1$.

\item
Suppose that $v \in V\1$. Then $\ip{0,v} = 0$, which implies that $v \in \{0\}^\perp\1$. Thus $\{0\}^\perp = V\1$.

\item
Suppose that $v \in V^\perp\1$. Then $\ip{v,v} = 0$, which implies that $v = 0$. Thus $V^\perp = \{0\}$.

\item
Suppose $U$ is a subset of $V$ and $u \in U \cap U^\perp\1$. Then $\ip{u,u} = 0$, which implies that $u = 0$. Thus $U \cap U^\perp \subset \{0\}$.

\item
Suppose $G$ and $H$ are subsets of $V$ and $G \subset H$. Suppose $v \in H^\perp\1$. Then $\ip{u,v} = 0$ for every $u \in H$, which implies that $\ip{u,v} = 0$ for every $u \in G$. Hence $v \in G^\perp\1$. Thus $H^\perp \subset G^\perp\1$.
\qed
\end{subtheorem}

Recall that if $U$ and $W$ are subspaces of $V\1$, then $V$ is the direct sum of $U$ and $W$ (written $V = U \oplus W$) if each element of $V$ can be written in exactly one way as a vector in $U$ plus a vector in $W$ (see \ref{directsum}). Furthermore, this happens if and only if $V = U + W$ and $U \cap W = \br{0}$ (see  \ref{219}).

The next result shows that every finite-dimensional subspace of $V$ leads to a natural direct sum decomposition of~$V\1$. See Exercise \ref{dense} for an example showing that the result below can fail without the hypothesis that the subspace $U$ is finite-dimensional.

\ResultBox[191]{direct sum of a subspace and its orthogonal complement}{\index{direct sum!of a subspace and its orthogonal complement}%
Suppose $U$ is a finite-dimensional subspace of $V\1$. Then
\[
V = U \oplus U^\perp\1.
\]
}

\proof
First we will show that
\[
V = U + U^\perp\1.
\]
To do this, suppose that $v \in V\1$.  Let $e_1, \dots, e_m$ be an orthonormal basis of~$U$. We want to write $v$ as the sum of a vector in $U$ and a vector orthogonal to $U$.

\pagebreak[3]

We have
\begin{equation} \label{416}
v = \underbrace{\ip{v, e_1}e_1 + \dots + \ip{v, e_m}e_m}_u
+ \underbrace{v - \ip{v, e_1}e_1 - \dots - \ip{v, e_m}e_m}_w\!.
\end{equation}%
Let $u$ and $w$ be defined as in the equation above (as was done in the proof of \ref{Bessel}). Because each $e_k \in U$, we see that $u \in U$.  Because $e_1, \dots, e_m$ is an orthonormal
list, for each $k = 1, \dots, m$ we have
\begin{align*}
\ip{w, e_k} &= \ip{v, e_k} - \ip{v, e_k} \\[3bp]
&= 0.
\end{align*}
Thus $w$ is orthogonal to every vector in $\Span(e_1, \dots, e_m)$, which shows that $w \in U^\perp\1$. Hence we have written
$v = u + w$, where $u \in U$ and $w \in U^\perp\1$, completing the proof that $V = U + U^\perp\1$.

From \ref{oc}(d), we know that $U \cap U^\perp = \{0\}$. Now equation $V = U + U^\perp\1$ implies that $V = U \oplus U^\perp$ (see \ref{219}).
\qed

Now we can see how to compute $\dim U^\perp$ from $\dim U$.

\ResultBox[350]{dimension of orthogonal complement}
{Suppose $V$ is finite-dimensional and $U$ is a subspace of $V\1$. Then
\[
\dim U^\perp = \dim V - \dim U.
\]}

\proof
The formula for $\dim U^\perp$ follows immediately from \ref{191} and \ref{3directsumdim}. \qed

The next result is an important consequence of \ref{191}.

\ResultBox[376]{orthogonal complement of the orthogonal complement}{%
Suppose $U$ is a finite-dimensional subspace of $V\1$. Then
\[
U = \paren[\big]{U^\perp}^\perp\1.
\]
}

\proof
First we will show that
\begin{equation} \label{415}
U \subset \paren[\big]{U^\perp}^\perp\1.
\end{equation}%
To do this, suppose $u \in U$.  Then $\ip{u, w} = 0$ for every $w \in U^\perp$ \big(by the definition of $U^\perp$\big).  Because $u$ is
orthogonal to every vector in $U^\perp\1$, we have $u \in \paren[\big]{U^\perp}^\perp\1$, completing the proof of \ref{415}.

To prove the inclusion in the other direction, suppose $v \in \paren[\big]{U^\perp}^\perp\1$.  By \ref{191}, we can write $v = u + w$, where
$u \in U$ and $w \in U^\perp\1$.  We have $v - u = w \in U^\perp\1$.  Because
$v \in \paren[\big]{U^\perp}^\perp$ and $u \in \paren[\big]{U^\perp}^\perp$
(from~\ref{415}), we have $v - u \in \paren[\big]{U^\perp}^\perp\1$.  Thus
$v - u \in  U^\perp \cap \paren[\big]{U^\perp}^\perp\1$,
which implies that $v - u = 0$ [by \ref{oc}(d)], which implies that $v = u$, which implies that
$v \in  U$.  Thus $\paren[\big]{U^\perp}^\perp \subset U$, which along with \ref{415}
completes the proof.
\qed

\pagebreak[3]

\mar{Exercise \ref{dense}\uppar{a} shows that the result below is not true without the hypothesis that $U$ is finite-dimensional.}

Suppose $U$ is a subspace of $V$ and we want to show that $U = V\1$. In some situations, the easiest way to do this is to show that the only vector orthogonal to $U$ is $0$, and then use the result below. For example, the result below is useful for Exercise \ref{190converse}.

\ResultBox[01Uperp]{$U^\perp = \{0\} \iff U = V$ \uppar{for\/ $U$ a finite-dimensional subspace of\/ $V$}}
{Suppose $U$ is a finite-dimensional subspace of $V\1$.  Then
\[
U^\perp = \{0\} \iff U = V\1.
\]}

\proof
First suppose $U^\perp = \br{0}$. Then by \ref{376}, $U = \paren[\big]{U^\perp}^\perp = \br{0}^\perp= V\1$, as desired.

Conversely, if $U = V\1$, then $U^\perp = V^\perp = \br{0}$ by \ref{oc}(c).
\qed

We now define an operator $P_{\!U}$ for each finite-dimensional subspace $U$ of $V\1$.

\DefinitionBox[orthogproj]{orthogonal projection, $P_{\!U}$\sindex{PU@$P_{\hspace{-.16667em}U}$}}
{Suppose $U$ is a finite-dimensional subspace of $V\1$. The \emph{orthogonal \mbox{projection}}\index{orthogonal!projection} of~$V$ onto~$U$ is the operator $P_{\!U} \in \L(V)$ defined as follows: For each $v \in V\1$, write $v = u + w$, where $u \in U$ and $w \in U^\perp\1$. Then let $P_{\!U} v = u$.}

The direct sum decomposition $V = U \oplus U^\perp$ given by \ref{191} shows that each $v \in V$ can be uniquely written in the form
$v = u + w$ with $u \in U$ and $w \in U^\perp\1$. Thus $P_{\!U} v$ is well defined. See the figure that accompanies the proof of \ref{417} for the picture describing $P_{\!U} v$ that you should keep in mind.

\begin{example}{orthogonal projection onto one-dimensional subspace}
Suppose $u \in V$ with $u \ne 0$ and $U$ is the one-dimensional subspace of $V$ defined by $U = \Span(u)$.

If $v \in V\1$, then
\[
v = \frac{\ip{v, u}}{\|u\|^2} u + \paren[\bigg]{ v - \frac{\ip{v, u}}{\|u\|^2} u },
\]
where the first term on the right is in $\Span(u)$ (and thus is in $U$) and the second term on the right is orthogonal to $u$ (and thus is in $U^\perp)$. Thus $P_{\!U} v$ equals the first term on the right. In other words, we have the formula
\[
P_{\!U} v = \frac{\ip{v, u}}{\|u\|^2} u
\]
for every $v \in V\1$.

The formula above becomes $P_{\!U} u = u$ if $v = u$ and becomes $P_{\!U} v = 0$ if $v \in \br{u}^\perp\1$. These equations are special cases of (b) and (c) in the next result.\looseness=-1
\end{example}

\pagebreak[3]

\ResultBox[418]{properties of orthogonal projection $P_{\!U}$}
{Suppose $U$ is a finite-dimensional subspace of $V\1$. Then
\begin{subtheorem}* \addtolength{\itemsep}{-.15bp} %VisualFormatting for page break

\item
$P_{\!U} \in \L(V)$;

\item
$P_{\!U} u = u$ for every $u \in U$;

\item
$P_{\!U} w = 0$ for every $w \in U^\perp$;

\item
$\ran P_{\!U} = U$;

\item
$\ker P_{\!U} = U^\perp$;

\item
$v - P_{\!U} v \in U^\perp$ for every $v \in V$;

\item
${P_{\!U}}^{\2} = P_{\!U}$;

\item
$\|P_{\!U} v\| \le \|v\|$ for every $v \in V$;

\item
if $e_1, \dots, e_m$ is an orthonormal basis of $U$ and $v \in V\1$, then
\[
P_{\!U} v = \ip{v, e_1} e_1 + \dots + \ip{v, e_m} e_m.
\]
\end{subtheorem}}

\pagebreak[3]

\proof
\begin{subtheorem}*
\item
To show that $P_{\!U}$ is a linear map on $V\1$, suppose $v_1, v_2 \in V\1$. Write
\[
v_1 = u_1 + w_1 \andd v_2 = u_2 + w_2
\]
with $u_1, u_2 \in U$ and $w_1, w_2 \in U^\perp\1$. Thus $P_{\!U} v_1 = u_1$ and $P_{\!U} v_2 = u_2$. Now
\[
v_1 + v_2 = (u_1 + u_2) + (w_1 + w_2),
\]
where $u_1 + u_2 \in U$ and $w_1 + w_2 \in U^\perp\1$. Thus
\[
P_{\!U}(v_1 + v_2) = u_1 + u_2 = P_{\!U} v_1 + P_{\!U} v_2.
\]

Similarly, suppose $\la \in \F{}$ and $v \in V\1$. Write $v = u + w$, where $u \in U$ and $w \in U^\perp\1$. Then $\la v = \la u + \la w$ with $\la u \in U$ and $\la w \in U^\perp\1$. Thus $P_{\!U}(\la v) = \la u = \la P_{\!U}v$.

Hence $P_{\!U}$ is a linear map from $V$ to $V\1$.

\item
Suppose $u \in U$. We can write $u = u + 0$, where $u \in U$ and $0 \in U^\perp\1$. Thus $P_{\!U}u = u$.

\item
Suppose $w \in U^\perp\1$. We can write $w = 0 + w$, where $0 \in U$ and $w \in U^\perp\1$. Thus $P_{\!U}w = 0$.

\item
The definition of $P_{\!U}$ implies that $\ran P_{\!U} \subset U$. Furthermore, (b) implies that $U \subset \ran P_{\!U}$. Thus $\ran P_{\!U} = U$.

\item
The inclusion $U^\perp \subset \ker P_{\!U}$ follows from (c). To prove the inclusion in the other direction, note that if $v \in \ker P_{\!U}$ then the decomposition given by \ref{191} must be $v = 0 + v$, where $0 \in U$ and $v \in U^\perp\1$. Thus $\ker P_{\!U} \subset U^\perp\1$.

\item
If $v \in V$ and $v = u + w$ with $u \in U$ and $w \in U^\perp\1$, then
\[
v - P_{\!U}v = v - u = w \in U^\perp\1.
\]

\item
If $v \in V$ and $v = u + w$ with $u \in U$ and $w \in U^\perp\1$, then
\[
\paren[\big]{{P_{\!U}}^{\2}}v = P_{\!U}(P_{\!U} v) = P_{\!U} u = u = P_{\!U}v.
\]

\item
If $v \in V$ and $v = u + w$ with $u \in U$ and $w \in U^\perp\1$, then
\[
\|P_{\!U} v\|^2 = \|u\|^2 \le \|u\|^2 + \|w\|^2 = \|v\|^2\1,
\]
where the last equality comes from the Pythagorean theorem.

\item
The formula for $P_{\!U} v$ follows from equation \ref{416} in the proof of \ref{191}.
\qed
\end{subtheorem}

In the previous section we proved the Riesz representation theorem (\ref{50}), whose key part states that every linear functional on a finite-dimensional inner product space is given by taking the inner product with some fixed vector. Seeing a different proof often provides new insight. Thus we now give a new proof of the key part of the Riesz representation theorem using orthogonal complements instead of orthonormal bases as in our previous proof.

The restatement below of the Riesz representation theorem provides an identification of $V$ with $V'\!$. We will prove only the ``onto'' part of the result below because the more routine ``one-to-one'' part of the result can be proved as in \ref{50}.

Intuition behind this new proof: If $\phi \in V'\!$, $v \in V\1$, and $\phi(u) = \ip{u,v}$ for all $u \in V\1$, then
$v \in (\ker \phi)^\perp\1$. However, $(\ker \phi)^\perp$ is a one-dimensional subspace of $V$ (except for the trivial case in which $\phi = 0$), as follows from \ref{350} and \ref{3ab11}. Thus we can obtain $v$ by choosing any nonzero element of $(\ker \phi)^\perp$ and then multiplying by an appropriate scalar, as is done in the proof below.

\ResultBox[3RRT]{Riesz representation theorem, revisited}
{Suppose $V$ is finite-dimensional. For each $v \in V\1$, define $\phi_v \in V'$ by\index{Riesz representation theorem}
\[
\phi_v(u) = \ip{u, v}
\]
for each $u \in V\1$. Then $v \mapsto \phi_v$ is a one-to-one function from $V$ onto $V'\!$.}

\mar{\textbf{Caution:} The function\/ $v \mapsto \phi_v$ is a linear mapping from\/ $V$ to\/ $V'$ if\/ \mbox{$\F{} = \R{}$}. However, this function is not linear if\/ $\F{}=\C{}$ because\/ $\phi_{\la v} = \bar{\la} \phi_v$ if\/ $\la \in \C{}$.}

\proof
To show that $v \mapsto \phi_v$ is surjective, suppose $\phi \in V'\!$.
If $\phi = 0$, then $\phi = \phi_0$. Thus assume $\phi \ne 0$. Hence $\ker \phi \ne V\1$, which implies that
$(\ker \phi)^\perp \ne \br{0}$ (by \ref{191} with $U = \ker \phi$).

Let $w \in (\ker \phi)^\perp$ be such that $w \ne 0$. Let
\begin{equation} \label{wu}
v = \frac{ \overline{\phi(w)} }{\norm{w}^2} w.
\end{equation}%
Then $v \in (\ker \phi)^\perp\1$. Also, $v \ne 0$ (because $w \notin \ker \phi$).

\pagebreak[3]

Taking the norm of both sides of \ref{wu} gives
\begin{equation} \label{normw}
\norm{v} = \frac{ \abs{\phi(w) }}{\norm{w}}.
\end{equation}%
Applying $\phi$ to both sides of \ref{wu} and then using \ref{normw}, we have
\[
\phi(v) = \frac{ \abs{\phi(w)}^2 }{ \norm{w}^2} = \norm{v}^2\1.
\]

Now suppose $u \in V\1$. Using the equation above, we have
\[
u = \paren[\bigg]{ u - \frac{ \phi(u) }{ \phi(v) } v } + \frac{ \phi(u) }{ \norm{v}^2 } v.
\]
The first term in parentheses above is in $\ker \phi$ and hence is orthogonal to $v$. Thus taking the inner product of both sides of the equation above with $v$ shows that
\[
\ip{u, v} = \frac{ \phi(u) }{ \norm{v}^2 } \ip{v, v} = \phi(u).
\]
Thus $\phi = \phi_v$, showing that $v \mapsto \phi_v$ is surjective, as desired.
\qed

See Exercise~\ref{VVpiso} for yet another proof of the Riesz representation theorem.

\pagebreak[3]

\subsection{Minimization Problems}

\mar[6]{The remarkable simplicity of the solution to this minimization problem has
led to many important applications of inner product spaces outside of pure
mathematics.}

The following problem often arises: Given a subspace $U$ of $V$ and a point
$v \in V\1$, find a point $u \in U$ such that $\|v - u\|$ is as small as
possible. The next result shows that $u = P_{\!U} v$ is the unique solution of this minimization
problem.\index{minimizing distance}

\ResultBox[417]{minimizing distance to a subspace}{%
Suppose\/ $U$ is a finite-dimensional subspace of\/ $V\1$, $v \in V\1$, and $u \in U$.  Then
\[
\| v - P_{\!U} v \| \le \| v - u \|.
\]
Furthermore, the inequality above is an equality if and only if $u = P_{\!U} v$.}

\proof
We have
\begin{align}
\| v - P_{\!U} v \|^2 &\le \|v - P_{\!U} v\|^2 + \|P_{\!U} v - u\|^2 \label{422}\\[5bp]
& = \norm[\big]{(v - P_{\!U} v) + (P_{\!U} v - u)}^2 \notag\\[5bp]
&= \| v - u \|^2\1, \notag
\end{align}
\marF{\includegraphics{Mathematica/ClosestPoint.pdf}}{$P_{\!U} v$ is the closest\\point in\/ $U$ to\/ $v$.}
where the first line above holds because $0 \le \|P_{\!U} v - u\|^2\1$, the second line above comes from the Pythagorean theorem \sbr{which applies because $v - P_{\!U} v \in U^\perp$ by \ref{418}(f), and $P_{\!U} v - u \in U$}, and the third line above holds by simple algebra.  Taking square roots gives the desired inequality.

The  inequality proved above is an equality if and only if \ref{422} is an equality, which
happens if and only if \mbox{$\|P_{\!U} v - u\| = 0$}, which happens if and only
if \mbox{$u = P_{\!U} v$}.
\qed

The last result is often combined with the formula \ref{418}(i) to
compute explicit solutions to minimization problems, as in the following example.

\begin{example}{using linear algebra to approximate the sine function}
Suppose we want to find a polynomial $u$ with real coefficients and of degree at most $5$ that approximates the sine function as well as possible on the interval $[-\pi, \pi]$, in the sense that
\[
\int_{-\pi}^{\pi} \abs[\big]{\sin x - u(x)}^2\,dx
\]
is as small as possible.

Let $C[-\pi, \pi]$ denote the real inner product space of continuous real-valued functions on $[-\pi, \pi]$ with inner product
\begin{equation} \label{419}
\ip{f, g} = \int_{-\pi}^{\pi} fg.
\end{equation}%
Let $v \in C[-\pi, \pi]$ be the function defined by $v(x) = \sin x$.
Let $U$ denote the subspace of $C[-\pi, \pi]$ consisting of the polynomials with real coefficients and of degree at most~$5$.  Our problem can now be reformulated as follows:
\[
\text{Find $u\in U$ such that $\|v - u\|$ is as
small as possible.}
\]

\mar{A computer that can integrate is useful here.}

To compute the solution to our approximation problem, first apply the
Gram--Schmidt procedure (using the inner product given by \ref{419}) to
the basis $1, x, x^2\1, x^3\1, x^4\1, x^5$ of $U$, producing an orthonormal basis
$e_1, e_2, e_3, e_4, e_5, e_6$ of $U$. Then, again using the inner
product given by~\ref{419}, compute $P_{\!U} v$ using \ref{418}(i) (with
$m = 6$).  Doing this computation shows that $P_{\!U} v$ is  the function $u$ defined by
\begin{equation} \label{420}
u(x) = 0.987862 x - 0.155271 x^3 + 0.00564312 x^5\1,
\end{equation}%
where the $\pi$'s that appear in the exact answer have been replaced
with a good decimal approximation.
By \ref{417}, the polynomial $u$ above is the best approximation to the sine function on $[-\pi, \pi]$ using
polynomials of degree at most~$5$ \big(here ``best approximation'' means in the sense of minimizing
$\int_{-\pi}^{\pi} |\sin x - u(x)|^2\,dx$\big).

To see how good this approximation is, the next figure shows the graphs of both the sine function and our approximation $u$ given by \ref{420} over the interval $[-\pi, \pi]$.

\FigureHereUnround{Mathematica/SinGram.pdf}{1}{Graphs on\/ $[-\pi, \pi]$ of the sine function \uppar{red} and its best\\fifth degree polynomial approximation\/ $u$ \uppar{blue} from~\ref{420}.}

Our approximation \ref{420} is so accurate that the two
graphs are almost identical---our eyes may see only one graph! Here the red graph is placed almost exactly over the blue graph. If you are viewing this on an electronic device, enlarge the picture above by 400\% near $\pi$ or $-\pi$ to see a small gap between the two graphs.

Another well-known approximation to the sine function by a polynomial of degree $5$ is given by the Taylor  polynomial $p$ defined by
\begin{equation} \label{421}
p(x) = x - \frac{x^3}{3!}  + \frac{x^5}{5!}.
\end{equation}%
To see how good this approximation is, the next picture shows the graphs of both the sine function and the Taylor
polynomial $p$ over the interval $[-\pi, \pi]$.

\FigureHereUnround{Mathematica/SinTaylor.pdf}{1}{Graphs on\/ $[-\pi, \pi]$ of the sine function \uppar{red}\\and the Taylor polynomial \uppar{blue} from \ref{421}.}

The Taylor polynomial of degree $5$ is an excellent approximation to $\sin x$ for $x$ near~$0$.  But the picture above shows that for $|x| > 2$, the Taylor polynomial is not so accurate, especially compared to \ref{420}.  For example, taking $x = 3$, our approximation \ref{420} estimates $\sin 3$ with an error of approximately $0.001$, but the Taylor polynomial \ref{421} estimates $\sin 3$ with an error of approximately $0.4$. Thus at $x = 3$, the error in the Taylor polynomial is hundreds of times larger than the error given by \ref{420}.  Linear algebra has helped us discover an approximation to the sine function that improves upon what we learned in calculus!
\end{example}

\subsection{Pseudoinverse}

Suppose $T \in \L(V\1, W)$ and $w \in W\3$. Consider the problem of finding $v \in V$ such that
\[
Tv = w.
\]
For example, if $V = \F n$ and $W = \FF m$, then the equation above could represent a system of $m$ linear equations in $n$ unknowns $v_1, \ldots, v_n$, where $v = (v_1, \ldots, v_n)$.

If $T$ is invertible, then the unique solution to the equation above is $v = T^{-1}w$. However, if $T$ is not invertible, then for some
$w \in W$ there may not exist any solutions of the equation above, and for some $w \in W$ there may exist infinitely many solutions of the equation above.

If $T$ is not invertible, then we can still try to do as well as possible with the equation above. For example, if the equation above has no solutions, then instead of solving the equation $Tv - w = 0$, we can try to find $v \in V$ such that $\norm{Tv - w}$ is as small as possible. As another example, if the equation above has infinitely many solutions $v \in V\1$, then among all those solutions we can try to find one such that $\norm{v}$ is as small as possible.

The pseudoinverse will provide the tool to solve the equation above as well as possible, even when $T$ is not invertible. We need the next result to define the pseudoinverse.

In the next two proofs, we will use without further comment the result that if $V$ is finite-dimensional and $T \in \L(V\1,W)$, then $\ker T$, $(\ker T)^\perp\1$, and $\ran T$ are all finite-dimensional.

\ResultBox{restriction of a linear map to obtain a one-to-one and onto map}
{Suppose $V$ is finite-dimensional and $T \in \L(V\1, W)$. Then $T|_{(\ker T)^\perp}$ is an injective map of $(\ker T)^\perp$ onto $\ran T\3$.}

\proof
Suppose that $v \in (\ker T)^\perp$ and $T|_{(\ker T)^\perp}v = 0$. Hence $Tv = 0$ and thus $v \in (\ker T) \cap (\ker T)^\perp\1$, which implies that
$v = 0$ \sbr{by \ref{oc}(d)}. Hence $\ker T|_{(\ker T)^\perp} = \br{0}$, which implies that $T|_{(\ker T)^\perp}$ is injective, as desired.

Clearly $\ran T|_{(\ker T)^\perp} \subset \ran T\3$. To prove the inclusion in the other direction, suppose $w \in \ran T\3$. Hence there exists $v \in V$ such that $w = Tv$. There exist $u \in \ker T$ and $x \in (\ker T)^\perp$ such that $v = u + x$ (by \ref{191}). Now
\[
T|_{(\ker T)^\perp}x = Tx = Tv - Tu = w - 0 = w,
\]
which shows that $w \in \ran T|_{(\ker T)^\perp}$. Hence $\ran T \subset \ran T|_{(\ker T)^\perp}$, completing the proof that
$\ran T|_{(\ker T)^\perp} = \ran T\3$.
\qed

\mar[3]{To produce the pseudoinverse notation $T^\dagger$ in \TeX, type \textup{\texttt{T\^\ \backslash dagger}}.}

Now we can define the pseudoinverse $T^\dagger$ (pronounced ``$T$ dagger'') of a linear map $T\3$. In the next definition (and from now on), think of $T|_{(\ker T)^\perp}$ as an invertible linear map from $(\ker T)^\perp$ onto $\ran T$, as is justified by the result above.

\DefinitionBox[6pseudo]{pseudoinverse, $T^\dagger$}
{Suppose that $V$ is finite-dimensional and $T \in \L(V\1, W)$. The \emph{pseudoinverse} $T^\dagger \in \L(W\3, V)$ of $T$ is the linear map from $W$ to $V$ defined by
\index{pseudoinverse}\sindex{Tdagger@$T^\dagger\1$}
\[
T^\dagger w = (T|_{(\ker T)^\perp})^{-1} P_{\!\ran T} \, w
\]
for each $w \in W\3$.}

Recall that $P_{\!\ran T} \, w = 0$ if $w \in (\ran T)^\perp$ and $P_{\!\ran T} \, w = w$ if $w \in \ran T\3$. Thus if $w \in (\ran T)^\perp\1$, then $T^\dagger w = 0$, and if $w \in \ran T$, then $T^\dagger w$ is the unique element of $(\ker T)^\perp$ such that
$T \paren[\big] {T^\dagger w} = w$.

The pseudoinverse behaves much like an inverse, as we will see.

\ResultBox[6propps]{algebraic properties of the pseudoinverse}
{Suppose $V$ is finite-dimensional and $T \in \L(V\1,W)$.
\begin{subtheorem}*
\item
If $T$ is invertible, then $T^\dagger = T^{-1}\1$.
\item
$TT^\dagger = P_{\!\ran T} =$ the orthogonal projection of $W$ onto $\ran T\3$.
\item
$T^\dagger T = P_{\!(\ker T)^\perp} =$ the orthogonal projection of $V$ onto $(\ker T)^\perp\1$.
\end{subtheorem}}

\proof
\begin{subtheorem} \addtolength{\itemsep}{-2bp}
\item
Suppose $T$ is invertible. Then $(\ker T)^\perp = V$ and $\ran T = W\3$. Thus $T|_{(\ker T)^\perp}= T$ and $P_{\!\ran T}$ is the identity operator on $W\3$. Hence $T^\dagger = T^{-1}\1$.

\item
Suppose $w \in \ran T\3$. Thus
\[
T T^\dagger w = T (T|_{(\ker T)^\perp})^{-1} w = w = P_{\!\ran T} \, w.
\]

If $w \in (\ran T)^\perp\1$, then $T^\dagger w = 0$. Hence $TT^\dagger w = 0 = P_{\!\ran T} \, w$.
Thus $TT^\dagger$ and $P_{\!\ran T}$ agree on $\ran T$ and on $(\ran T)^\perp\1$. Hence these two linear maps are equal (by \ref{191}).

\item
Suppose $v \in (\ker T)^\perp\1$. Because $Tv \in \ran T\3$, the definition of $T^\dagger$ shows that
\[
T^\dagger(Tv)= (T|_{(\ker T)^\perp})^{-1} (Tv) = v = P_{\!(\ker T)^\perp} v.
\]

If $v \in \ker T$, then $T^\dagger T v = 0 = P_{\!(\ker T)^\perp} v$.
Thus $T^\dagger T$ and $P_{\!(\ker T)^\perp}$ agree on $(\ker T)^\perp$ and on $\ker T\3$. Hence these two linear maps are equal (by \ref{191}).
\qed
\end{subtheorem}

\mar{The pseudoinverse is also called the Moore--Penrose inverse.\index{Moore--Penrose inverse}}

Suppose that $T \in \L(V\1, W)$. If $T$ is surjective, then $TT^\dagger$ is the identity operator on~$W\3$, as follows from (b) in the result above. If $T$ is injective, then $T^\dagger T$ is the identity operator on $V\1$, as follows from (c) in the result above. For additional algebraic properties of the pseudoinverse, see Exercises \ref{3dagger}--\ref{2dag}.

\pagebreak[3]

For $T \in \L(V\1,W)$ and $w \in W\3$, we now return to the problem of finding $v \in V$ that solves the equation
\[
Tv = w.
\]
As we noted earlier, if $T$ is invertible, then $v = T^{-1}w$ is the unique solution, but if $T$ is not invertible, then $T^{-1}$ is not defined. However, the pseudoinverse $T^\dagger$ is defined. Taking $v = T^\dagger w$ makes $Tv$ as close to $w$ as possible, as shown by (a) of the next result. Thus the pseudoinverse provides what is called a \textit{best fit} to the equation above.

Among all vectors $v \in V$ that make $Tv$ as close as possible to $w$, the vector $T^\dagger w$ has the smallest norm, as shown by combining (b) in the next result with the condition for equality in (a).

\ResultBox{pseudoinverse provides best approximate solution or best solution}
{Suppose $V$ is finite-dimensional, $T \in \L(V\1,W)$, and $w \in W\3$.
\begin{subtheorem}*
\item
If $v \in V\1$, then
\[
\norm[\big]{ T(T^\dagger w) - w } \le \norm{Tv - w},
\]
with equality if and only if $v \in T^\dagger w + \ker T\3$.
\item
If $v \in T^\dagger w + \ker T$, then
\[
\norm[\big]{T^\dagger w} \le \norm{v},
\]
with equality if and only if $v = T^\dagger w$.
\end{subtheorem}}

\proof
\begin{subtheorem}
\item
Suppose $v \in V\1$. Then
\[
Tv - w = (Tv - TT^\dagger w) + (TT^\dagger w - w).
\]
The first term in parentheses above is in $\ran T\3$. Because the operator $TT^\dagger$ is the orthogonal projection of $W$ onto $\ran T$ \sbr{by \ref{6propps}(b)}, the second term in parentheses above is in $(\ran T)^\perp$ \sbr{see \ref{418}(f)}.

Thus the Pythagorean theorem implies the desired inequality that the norm of the second term in parentheses above is less than or equal to $\norm{Tv - w}$, with equality if and only if the first term in parentheses above equals $0$. Hence we have equality if and only if $v - T^\dagger w \in \ker T$, which is equivalent to the statement that $v \in T^\dagger w + \ker T$, completing the proof of (a).

\item
Suppose $v \in T^\dagger w + \ker T\3$. Hence $v - T^\dagger w \in \ker T\3$. Now
\[
v = (v - T^\dagger w) + T^\dagger w.
\]
The definition of $T^\dagger$ implies that $T^\dagger w \in (\ker T)^\perp\1$. Thus the Pythagorean theorem implies that
$\norm[\big]{T^\dagger w} \le \norm{v}$, with equality if and only if $v = T^\dagger w$.
\qed
\end{subtheorem}

A formula for $T^\dagger$ will be given in the next chapter (see \ref{7SVD}).

\begin{example}{pseudoinverse of a linear map from\/ $\F4$ to\/ $\F3$}
Suppose $T \in \L\paren[\big]{ \F4, \F3 }$ is defined by
\[
T(a, b, c, d) = (a + b + c, 2c + d, 0).
\]
This linear map is neither injective nor surjective, but we can compute its pseudo\-inverse. To do this, first note that
$\ran T = \br{(x,y,0) : x, y \in \F{}}$. Thus
\[
P_{\ran T}(x,y,z) = (x, y, 0)
\]
for each $(x, y, z) \in \FF3$.
Also,
\[
\ker T = \br[\big]{ (a, b, c, d) \in \F4 : a+b+c= 0 \text{ and } 2c + d = 0}.
\]
The list $(-1, 1, 0, 0), (-1, 0, 1, -2)$ of two vectors in $\ker T$ spans $\ker T$ because if $(a, b, c, d) \in \ker T$ then
\[
(a, b, c, d) = b(-1, 1, 0, 0) + c(-1,0,1,-2).
\]
Because the list $(-1, 1, 0, 0), (-1, 0, 1, -2)$ is linearly independent, this list is a basis of $\ker T\3$.

Now suppose $(x, y, z) \in \FF3$. Then
\begin{equation} \label{6Tpseudo}
T^\dagger(x, y, z) = (T|_{(\ker T)^\perp})^{-1} P_{\!\ran T} (x, y, z) = (T|_{(\ker T)^\perp})^{-1} (x, y, 0).
\end{equation}%
The right side of the equation above is the vector $(a, b, c, d) \in \F4$ such that $T(a, b, c, d) = (x, y, 0)$ and
$(a, b, c, d) \in (\ker T)^\perp\1$. In other words, $a$, $b$, $c$, $d$ must satisfy the following equations:
\begin{align*}
a + b + c &=x \\
2c + d &= y \\
-a + b &= 0 \\
-a + c - 2d &= 0,
\end{align*}
where the first two equations are equivalent to the equation $T(a, b, c, d) = (x, y, 0)$ and the last two equations come from the condition for $(a, b, c, d)$ to be orthogonal to each of the basis vectors $(-1, 1, 0, 0), (-1, 0, 1, -2)$ in this basis of $\ker T\3$. Thinking of $x$ and $y$ as constants and $a$, $b$, $c$, $d$ as unknowns, we can solve the system above of four equations in four unknowns, getting
\[
a = \tfrac{1}{11} (5x-2y),\ b = \tfrac{1}{11} (5x-2y),\ c = \tfrac{1}{11} (x+4y),\ d = \tfrac{1}{11} (-2x + 3y).
\]
Hence \ref{6Tpseudo} tells us that
\[
T^\dagger(x,y,z) = \tfrac{1}{11} (5x-2y, 5x-2y, x+4y, -2x +3y).
\]
The formula above for $T^\dagger$ shows that $TT^\dagger(x, y, z) = (x, y, 0)$ for all $(x, y, z) \in \FF3$, which illustrates the equation $TT^\dagger = P_{\!\ran T}$ from \ref{6propps}(b).
\end{example}

\begin{exercises}

\exercise
Suppose $v_1, \dots, v_m \in V\1$. Prove that
\[
\{v_1, \dots, v_m\}^\perp = \bigl( \Span(v_1, \dots, v_m) \bigr)^\perp\1.
\]
\exercisedisplayend


\exercise
Suppose $U$ is a subspace of $V$ with basis $u_1, \dots, u_m$ and\label{6onbasis}
\[
u_1, \dots, u_m, v_1, \dots, v_n
\]
is a basis of $V\1$. Prove that if the Gram--Schmidt procedure is applied to the basis of $V$ above, producing a list $e_1, \dots, e_m, f_1, \dots, f_n$, then $e_1, \dots, e_m$ is an orthonormal basis of $U$ and $f_1, \dots, f_n$ is an orthonormal basis of $U^\perp\1$.


\exercise
Suppose $U$ is the subspace of $\R4$ defined by
\[
U = \Span\bigl( (1,2,3,-4), (-5,4, 3, 2) \bigr).
\]
Find an orthonormal basis of $U$ and an orthonormal basis of $U^\perp\1$.


\exercise
Suppose $e_1, \ldots, e_n$ is a list of vectors in $V$ with $\norm{e_k} = 1$ for each $k=1, \ldots, n$ and \label{190converse}
\[
\|v\|^2 = \abs[\big]{ \ip{v, e_1} }^2  + \dots + \abs[\big]{ \ip{v, e_n} }^2
\]
for all $v \in V\1$. Prove that $e_1, \ldots, e_n$ is an orthonormal basis of $V\1$.
\exercisecomment{This exercise provides a converse to \ref{190}\uppar{b}.}


\exercise
Suppose that $V$ is finite-dimensional and $U$ is a subspace of $V\1$. Show that
$P_{\!U^\perp} = I - P_{\!U}$, where $I$ is the identity operator on $V\1$.


\exercise
Suppose $V$ is finite-dimensional and $T \in \L(V\1, W)$. Show that
\[
T = TP_{\!(\ker T)^\perp}= P_{\!\ran T} T\3.
\]
\exercisedisplayend


\exercise
Suppose that $X$ and $Y$ are finite-dimensional subspaces of $V\1$. Prove that $P_{\!X} P_{\!Y} = 0$ if and only if $\ip{x, y} = 0$ for all $x \in X$ and all $y \in Y$.


\exercise
Suppose $U$ is a finite-dimensional subspace of $V$ and $v \in V\1$. Define a linear functional $\fun \phi U \F{}$ by
\[
\phi(u) = \ip{u, v}
\]
for all $u \in U$. By the Riesz representation theorem (\ref{50}) as applied to the inner product space $U$, there exists a unique vector
$w \in U$ such that\index{Riesz representation theorem}
\[
\phi(u) = \ip{u, w}
\]
for all $u \in U$. Show that $w = P_{\!U} v$.


\exercise
Suppose $V$ is finite-dimensional. Suppose $P \in \L(V)$ is such that $P^2 = P$ and every vector in
$\ker P$ is orthogonal to every vector in $\ran P\1$. Prove that there exists a subspace $U$ of $V$ such that $P = P_{\!U}$.\label{6orth}


\exercise
Suppose $V$ is finite-dimensional and $P \in \L(V)$ is such that $P^2 = P$ and
\[
\|Pv\| \le \|v\|
\]
for every $v \in V\1$.  Prove that there exists a subspace $U$ of $V$ such that $P = P_{\!U}$.


\exercise
Suppose $T \in \L(V)$ and $U$ is a finite-dimensional subspace of~$V\1$.  Prove that\label{6invP}
\[
U \text{ is invariant under } T \iff P_{\!U} T P_{\!U} = T P_{\!U}.
\]
\exercisedisplayend


\exercise
Suppose $V$ is finite-dimensional, $T \in \L(V)$, and $U$ is a subspace of~$V\1$.  Prove that
\[
U \text{ and } U^\perp \text{ are both invariant under } T \iff P_{\!U} T = T P_{\!U}.
\]
\exercisedisplayend


\exercise
Suppose $\F{} = \R{}$ and $V$ is finite-dimensional. For each $v \in V\1$, let $\phi_v$ denote the linear functional on $V$ defined by \label{VVpiso}\index{Riesz representation theorem}
\[
\phi_v(u) = \ip{u, v}
\]
for all $u \in V\1$.
\begin{subtheorem}
\item
Show that $v \mapsto \phi_v$ is an injective linear map from $V$ to~$V'\!$.
\item
Use (a) and a dimension-counting argument to show that $v \mapsto \phi_v$ is an isomorphism from $V$ onto~$V'\!$.
\end{subtheorem}
\exercisecomment{The purpose of this exercise is to give an alternative proof of the Riesz representation theorem \uppar{\ref{50} and \ref{3RRT}} when $\F{} = \R{}$. Thus you should not use the Riesz representation theorem as a tool in your solution.}


\exercise
Suppose that $e_1, \ldots, e_n$ is an orthonormal basis of $V\1$. Explain why the dual basis (see \ref{DualBasis}) of
$e_1, \ldots, e_n$ is $e_1, \ldots, e_n$ under the identification of $V'$ with $V$ provided by the Riesz representation theorem (\ref{3RRT}).


\exercise
In $\R{4}$, let
\[
U = \Span \bigl( (1, 1, 0, 0), (1, 1, 1, 2) \bigr).
\]
Find $u \in U$ such that $\| u - (1,2,3,4) \|$ is as small as possible.


\exercise
Suppose $C[-1, 1]$ is the vector space of continuous real-valued functions on the interval $[-1 ,1]$ with inner product given by \label{dense}
\[
\ip{f, g} = \int_{-1}^1 fg
\]
for all $f, g \in C[-1, 1]$. Let $U$ be the subspace of $C[-1, 1]$ defined by
\[ %\addtolength{\belowdisplayskip}{-3bp}
U = \br[\big]{ f \in C[-1, 1]: f(0) = 0 }.
\]
\begin{subtheorem}*
\item
Show that $U^\perp = \{0\}$.
\item
Show that \ref{191} and \ref{376} do not hold without the finite-dimensional hypothesis.
\end{subtheorem}
\exercisesubtheoremend


\exercise
Find $p \in \P\1_3(\R{})$ such that $p(0) = 0$, $p' \! (0) = 0$, and
$ \displaystyle
\int_0^1 \, \abs[\big]{2 + 3x - p(x)}^2 \,dx
$
is as small as possible.


\exercise
Find $p \in \P\1_5(\R{})$ that makes
$ \displaystyle
\int_{-\pi}^{\pi} \, \abs[\big]{\sin x - p(x)}^2\,dx
$
as small as possible.
\exercisecomment{The polynomial \ref{420} is an excellent
approximation to the answer to this exercise, but here you are asked
to find the exact solution, which involves powers of\/ $\pi$.  A
computer that can perform symbolic integration should help.}


\exercise
Suppose $V$ is finite-dimensional and $P \in \L(V)$ is an orthogonal projection of $V$ onto some subspace of $V\1$. Prove that $P^\dagger = P\1$.\label{3dagger}


\exercise
Suppose $V$ is finite-dimensional and $T \in \L(V\1, W)$. Show that \label{5NullRanTd}
\[
\ker T^\dagger = (\ran T)^\perp \andd \ran T^\dagger = (\ker T)^\perp\1.
\]
\exercisedisplayend


\exercise
Suppose $T \in \L\paren[\big]{ \F3, \F2 }$ is defined by
\[
T(a, b, c) = (a + b + c, 2b + 3c).
\]
\begin{subtheorem}*
\item
For $(x, y) \in \FF2$, find a formula for $T^\dagger(x,y)$.
\item
Verify that the equation $TT^\dagger = P_{\!\ran T}$ from \ref{6propps}(b) holds with the formula for $T^\dagger$ obtained in (a).
\item
Verify that the equation $T^\dagger T = P_{\!(\ker T)^\perp}$ from \ref{6propps}(c) holds with the formula for $T^\dagger$ obtained in (a).
\end{subtheorem}
\exercisesubtheoremend


\exercise
Suppose $V$ is finite-dimensional and $T \in \L(V\1, W)$. Prove that
\[
TT^\dagger T = T \andd T^\dagger T T^\dagger = T^\dagger\1.
\]
\exercisedisplayend
\exercisecomment{Both formulas above clearly hold if\/ $T$ is invertible because in that case we can replace\/ $T^\dagger$ with\/ $T^{-1}\1$.}


\exercise
Suppose $V$ and $W$ are finite-dimensional and $T \in \L(V\1, W)$. Prove that\label{2dag}
\[
\paren[\big]{T^\dagger}^\dagger = T\3.
\]
\exercisedisplayend
\exercisecomment{The equation above is analogous to the equation\/ $\paren[\big]{T^{-1}}^{-1} = T$ that holds if\/ $T$ is invertible.}


\end{exercises}

\CrCom%
\endinput
